% Standalone expansion of report.tex and its local inputs.
\documentclass[11pt]{article}
\usepackage[T1]{fontenc}
\usepackage{lmodern}
\usepackage[letterpaper,margin=0.9in]{geometry}
\usepackage{amsmath,amssymb,amsthm,mathtools}
\usepackage{booktabs,longtable,array,multicol}
\usepackage{microtype,xcolor,hyperref,url,fancyhdr}
\hypersetup{colorlinks=true,linkcolor=blue!45!black,urlcolor=blue!45!black,citecolor=blue!45!black,pdftitle={Quotient carriers and integral projection for Diophantine consistency},pdfauthor={Research prepared for Vladimir Reshetnikov with AI assistance}}
\pagestyle{fancy}\fancyhf{}\fancyhead[L]{\small Diophantine consistency}\fancyhead[R]{\small 25 October 2026}\fancyfoot[C]{\thepage}
\setlength{\headheight}{14pt}
\setlength{\parindent}{0pt}\setlength{\parskip}{5pt}
\setlength{\emergencystretch}{2em}
\allowdisplaybreaks[3]
\newtheorem{theorem}{Theorem}[section]
\newtheorem{lemma}[theorem]{Lemma}
\newtheorem{proposition}[theorem]{Proposition}
\theoremstyle{remark}\newtheorem{remark}[theorem]{Remark}
\newcommand{\Z}{\mathbb Z}\newcommand{\N}{\mathbb N}
\newcommand{\Con}{\operatorname{Con}}\newcommand{\lcm}{\operatorname{lcm}}
\newcommand{\Refutes}{\operatorname{Ref}}
\newcommand{\file}[1]{\nolinkurl{#1}}
\newcommand{\rg}[2]{\{#1,\ldots,#2\}}
\newcommand{\fl}[1]{\left\lfloor#1\right\rfloor}
% BEGIN results.tex
\newcommand{\DirectVars}{217}
\newcommand{\DirectLast}{216}
\newcommand{\DirectRows}{154}
\newcommand{\DirectRowLast}{153}
\newcommand{\DirectHeight}{5,184}
\newcommand{\DirectTerms}{19,112}

% END results.tex

\begin{document}
\hypersetup{pageanchor=false}
\begin{titlepage}
\vspace*{0.5in}
{\color{blue!45!black}\large Research contribution to ProveIt\par}
\vspace{0.25in}
{\LARGE\bfseries Quotient carriers and integral projection\par}
\vspace{0.12in}
{\Large Explicit polynomial equations for syntactic consistency of ZFC\par}
\vspace{0.35in}
{\large Prepared for Vladimir Reshetnikov\par}
{25 October 2026\par}
{Developed with AI assistance; ordinary proofs and exact finite artifacts\par}
\vspace{0.4in}
\begin{abstract}
We construct four compact ordinary polynomials and a direct integer
quartic whose unsolvability, in their stated domains, is provably
equivalent in ZFC to the syntactic consistency of ZFC. An auxiliary Pell
ordinate carries the residue of an omitted divisibility quotient. This
gives thirteen-positive-variable equations in 70 and 71 binary arithmetic
operations, of degrees 206 and 162. A smaller residue modulus gives
degree 154 in 72 operations. A separate rational carrier, made integral
by the inherited adjacent-Pell descent, gives twelve positive variables
in 74 operations at degree 438.

On the direct proof-calculus route, we combine an affine signed-radix
bootstrap with previously supplied strict local ranges, pooled power
certificates and integral quadratic projections. The resulting polynomial
has \DirectVars\ unrestricted integer variables, \DirectRows\ quadratic
residuals, degree four and coefficient height \DirectHeight. A separate
220-variable quartic has coefficient height 2,336. Neither direct
construction uses K215 or a universal Diophantine polynomial.

The report prints the complete principal residual list and the full
finite recipe for the compact fixed integers. Exact arithmetic programs,
all quartic coefficients, inverse maps, checkers and source identities
are supplied. We distinguish new applications from inherited lemmas,
finite verification from unbounded proofs, and ordinary formalizability
from a new proof-assistant derivation. No global minimum or consistency
decision is claimed.
\end{abstract}
\vfill
\begingroup\raggedright
\textbf{Comparison snapshot.}\quad
\file{536c08e7913e59bd42e8fca66e89944c920884ad}, including all twelve
archives present in \file{docs/incoming} at that commit.\par
\textbf{Canonical objects.}\quad
The literal equations, complete coefficient data and restoration
certificates in the accompanying archive.
\endgroup
\end{titlepage}
\pagenumbering{roman}
\hypersetup{pageanchor=true}
\begingroup\small\setlength{\parskip}{0pt}
\tableofcontents
\endgroup
\newpage
\pagenumbering{arabic}

\section{Results and the exact meaning of the theorem}

Let $\Refutes_{\mathrm{ZFC}}(p)$ be the fixed syntactic predicate expressing
that $p$ codes a finite derivation of contradiction in ZFC, with the
numeration and logical interfaces retained from ProveIt. Set
\[
 \Con(\mathrm{ZFC})\equiv
 \neg\exists p\in\N\;\Refutes_{\mathrm{ZFC}}(p).
\]
For each selected polynomial $P$ and its declared domain $D$, the result is
\begin{equation}\label{eq:target}
 \mathrm{ZFC}\vdash
 \left[\neg\exists\mathbf x\in D\;P(\mathbf x)=0\right]
 \longleftrightarrow\Con(\mathrm{ZFC}).
\end{equation}
We give constructions in both directions between a polynomial zero and
a finite refutation, by ordinary arithmetic formalizable in ZFC.
Neither direction assumes consistency. The theorem does not assert
that either side of the biconditional holds.

% BEGIN comparison.tex
\begin{center}\small
\begin{tabular}{@{}lrrr@{}}
\toprule
Compact equation & Positive variables & Operations & Degree\\
\midrule
Package 142, raw parent &14&70&189\\
Package 142, reduced parent &14&71&157\\
Package 142, anchored parent &13&72&433\\
\midrule
New residue carrier $P_{13,70}$ &13&70&206\\
New residue carrier $P_{13,71}$ &13&71&162\\
New smaller-modulus $P_{13,72}$ &13&72&154\\
New rational carrier $P_{12,74}$ &12&74&438\\
\bottomrule
\end{tabular}
\end{center}
All seven rows in this table use the same complete indexed943/radix67
constant tuple. The first two new forms remove one coordinate without
adding an arithmetic gate. At thirteen variables and 72 gates, the
smaller-modulus equation reduces the matched anchored parent's degree
from 433 to 154. The twelve-variable form uses two additional gates
relative to that parent. These are exact circuit properties, not
estimates from the appearance of a displayed formula.

\begin{center}\small
\begin{tabular}{@{}lrrrr@{}}
\toprule
Direct integer quartic & Variables & Residuals & Monomials & Height\\
\midrule
Incoming Pooled Powers &225&162&7,867&128,446\\
Incoming Quadratic Projection &234&171&19,050&4,480\\
\midrule
New principal $P_{217}$ &217&154&19,112&5,184\\
New $P_{220}$, lower height &220&157&8,847&2,336\\
New $P_{220}$, smaller support &220&157&7,858&6,272\\
\bottomrule
\end{tabular}
\end{center}
The quartic heights are maximum absolute coefficients after complete
expansion and collection. They are not residual heights. All coordinates
are unrestricted integers, and all these polynomials have degree four.
The 217-variable result trades expanded support for fewer variables and
rows. The smaller-support 220-variable result improves all four displayed
size measures over its actual 225-variable parent. The lower-height
220-variable result makes a different support/height tradeoff.

The direct results are new compositions of existing arithmetic and
projection methods with the specified source changes. In particular the
incoming Linear CRT archive already uses $F=t^t$, $R=F$, and a linear
odd beta multiplier. We do not claim priority for those choices.

% END comparison.tex


There is no single ordering of \emph{simplest} across variable count,
domain, degree, coefficient height, expanded support, and arithmetic
circuit size. The repository already contains a quartic with 37
unrestricted integer variables using a large fixed coefficient
architecture, and equations with nine positive variables at much
higher circuit cost and degree. The present direct quartics concern
the small-coefficient $C_{222}$ branch. The compact contribution improves
specific operation/degree tradeoffs and eliminates a variable from its
matched parents; it does not establish a global minimum.

All operation counts charge each binary addition, subtraction and
multiplication, including a square and multiplication by a small
constant. Reuse and fixed integer leaves are free. Thus the compact
counts do not charge the bit lengths of their large fixed constants.
Every comparison retains the variable domain and coefficient architecture.
The source ledger records actual pinned files, including incoming
manuscripts absent from older headline README tables.

\section{Compact equations by eliminating an index quotient}\label{sec:quotient}
% BEGIN compact_proof.tex
The compact contribution consists of four new descendants of the complete
indexed943/radix67 equations in package~142 \cite{coprime}. The inherited
Pell identities, coprime square extraction, and rational adjacent-Pell
descent are prior work. The new steps are the particular residue and
rational carriers which remove an index quotient, and the smaller modulus
which reduces the degree of one of the residue constructions.

\subsection{The complete polynomial definitions}

The four fixed positive integers $\beta,K_0,M_C,M_F$ are specified in
Appendix~\ref{app:compact} and supplied as complete finite byte strings.
They satisfy
\[
 \beta\ge31,\qquad K_0>0,\qquad 0<M_C<\beta<M_F<2\beta.
\]
These inequalities suffice for some early estimates. The full calibrated
tuple, rather than these inequalities alone, is required for the theorem.
For the first three equations the independent coordinates, all strictly
positive integers, are
\[
 (J,F,\alpha,z,f,n,i,T,s,w,\eta,\zeta,\gamma).
\]
Define the following polynomial abbreviations:
\begin{align*}
 q&=\beta J+1,& C&=q-F-\alpha,& X&=wq,&Y&=sq^3,\\
 E&=XY,&k&=\eta+\zeta,&\ell&=EkY,&c&=kY+\eta,\\
 a&=E+Y,&a_0&=a+2,&H&=4a+3,&\Delta&=a^2+H=a_0^2-1,\\
 D&=X+ac+\gamma H,&A&=1+\ell(\ell+k),
 \end{align*}
and
\begin{align*}
 R&=\bigl[(q-1)(q-F)+\alpha\bigr](q^2-1)+(M_C+qM_F)J,\\
 N_1&=D^2-\Delta c^2,\\
 N_t&=(K_0+w)C+(q-F)-z(q-1).
\end{align*}
The parent has an additional positive coordinate $h$ and factor
$N_k=k-hE-R$. Our new equations omit both. Every letter introduced by a
definition is an abbreviation, not a further quantified unknown.

\paragraph{The 70-operation equation.}
Put
\begin{align*}
 t&=En-k+R,& I&=Ei c^2,\\
 W&=\Delta A I^2,&Q&=\Delta W,&N_s&=f^2-W,\\
 V&=f^2(cT-R)-c,&B&=Q(t^2-AV^2)-t^2+A.
\end{align*}
Then
\begin{equation}\label{eq:p13raw}
 \boxed{P_{13,70}=N_1N_t(3B+1)N_s-1=0.}
\end{equation}
This is a polynomial in exactly thirteen positive unknowns, of exact
total degree 206. Its literal straight-line program uses 70 binary
additions, subtractions and multiplications.

\paragraph{The 71-operation equation.}
Retain $t=En-k+R$, and instead put
\begin{align*}
 m&=AEi,& I&=mc^2,& W&=\Delta I^2,\\
 N_s&=f^2-W,& Q&=\Delta(f^2-1),\\
 V&=f^2(cT-R)-c,&B&=Q(t^2-AV^2)-t^2+A.
\end{align*}
With these definitions use the same expression
\begin{equation}\label{eq:p13reduced}
 \boxed{P_{13,71}=N_1N_t(3B+1)N_s-1=0.}
\end{equation}
Its exact degree is 162, and its complete program has 71 operations.

\paragraph{The smaller-modulus equation.}
In \eqref{eq:p13reduced}, make just the two substitutions
\begin{equation}\label{eq:smallmod}
 t=q^4n-k+R,\qquad m=Aq^2i.
\end{equation}
Call the resulting polynomial $P_{13,72}$. Its exact degree is 154,
and its complete program has 72 operations. The extra operation computes
$q^4$ from the already present $q^2$.

\paragraph{The twelve-variable equation.}
Drop $f$ from the independent coordinate list. Keep all common outer
definitions, and define
\begin{align*}
 m&=a_0Di,& I&=mc^2,&W&=\Delta A I^2,\\
 F_2&=1+W,&Q&=\Delta W,&G&=AF_2,\\
 V&=F_2(cT-R)-c,&e&=k-R-1,&Z&=En+e,\\
 B_{12}&=Q(Z^2-E^2GV^2)-Z^2+E^2G.
\end{align*}
The equation is
\begin{equation}\label{eq:p12}
 \boxed{P_{12,74}=N_1N_t(3B_{12}+1)-1=0.}
\end{equation}
It has twelve positive unknowns, 74 operations and exact degree 438.
There is no rational operation in this definition. Rational quantities
will appear only in its proof.

\subsection{Recovering the first three Pell pairs}

At a zero of any of the three thirteen-variable equations, each factor
of the product in \eqref{eq:p13raw} or \eqref{eq:p13reduced} is an integer
unit. The factor $3B+1$ is 1 modulo 3, so it is $+1$ and $B=0$.
Furthermore $\Delta=a_0^2-1$ is 0 or 3 modulo 4. A norm
$u^2-\Delta v^2$ cannot be $-1$ modulo 4. Thus $N_1=1$.

The exact folded identity is
\begin{equation}\label{eq:fold}
 (Q-1)t^2=A(QV^2-1).
\end{equation}
In the raw construction, $A\mid Q$ and $Q>1$ hold unconditionally.
Equation~\eqref{eq:fold} excludes $t=0$. Both $Q-1$ and $QV^2-1$
are $-1$ modulo $A$. At each prime dividing $A$, valuation in
\eqref{eq:fold} therefore gives
$2v_p(t)=v_p(A)$. Consequently $A=\tau^2$ with $\tau>0$ integral,
and $\tau\mid t$. Only now may we write
$W=\Delta(\tau I)^2$ and exclude $N_s=-1$ modulo 4. Hence $N_s=1$.

In the reduced constructions, $W=\Delta I^2$ from the beginning, so
$N_s=1$ follows first. It gives $Q=(\Delta I)^2>1$. As $A\mid I$,
the same valuation argument again recovers $A=\tau^2$ and $\tau\mid t$.
In all cases $N_t=1$, since it is the only remaining factor.

Put $y=|t|/\tau$. With the positive strong ordinate $I_*$ and $S$ defined by
\[
 I_*=
 \begin{cases}
   \tau Ei c^2,&P_{13,70},\\
   AEi c^2,&P_{13,71},\\
   Aq^2i c^2,&P_{13,72},
 \end{cases}
 \qquad S=\Delta I_*,
\]
we have recovered
\begin{align}
 \tau^2&=1+\ell(\ell+k),& D^2-\Delta c^2&=1,\label{eq:firstmain}\\
 f^2-\Delta I_*^2&=1,&Q&=S^2,\label{eq:strong}\\
 QV^2-(Q-1)y^2&=1,&y&>0.\label{eq:aux}
\end{align}
Both signs of $t$ and $V$ have been allowed. Also $c^2\mid I_*$.

Before interpreting any index, $N_t=1$ gives
\[
 (K_0+w)C=1+z(q-1)-(q-F)\ge F>0.
\]
Thus $C>0$. Set $G_0=(q-1)(q-F)+\alpha=q\alpha+(q-1)C$.
The fixed constant bounds give
\[
 q\le G_0\le q(q-1),\qquad
 0<(M_C+qM_F)J<(2q+1)(q-1),
\]
and therefore
\begin{equation}\label{eq:Rbounds}
 3q<R<q^4-q^3+q^2-1<q^4-2\le E-2.
\end{equation}
In particular none of these bounds has used the omitted $h$.

\subsection{Authenticating the main index without that quotient}

For $b\ge2$, define the Pell sequences by
\[
 \chi_b(0)=1,\quad\chi_b(1)=b,\qquad
 \psi_b(0)=0,\quad\psi_b(1)=1,
 \qquad u_{r+2}=2bu_{r+1}-u_r.
\]
The standard solution classification, addition formulas, monotonicity,
and strong divisibility of $\psi_b$ are retained ordinary lemmas.
For clarity, the particular signed-index argument needed here is included.

Write $v=EY=XY^2$ and $P_0=2v+1$. The first equation in
\eqref{eq:firstmain} is $\tau^2-v(v+1)k^2=1$. Its fundamental
positive solution is $(2v+1,2)$: an ordinate of one would require a
square strictly between $v^2$ and $(v+1)^2$. Hence, for some $r\ge1$,
\[
 \tau=\chi_{P_0}(r),\qquad k=2\psi_{P_0}(r).
\]
The main norm gives $D=\chi_{a_0}(p)$ and $c=\psi_{a_0}(p)$ for
some $p\ge1$. If $r=1$, then $k=2$, $\eta=\zeta=1$, and
$c=2Y+1$ lies strictly between $\psi_{a_0}(1)=1$ and
$\psi_{a_0}(2)=2a_0$. This is impossible. Therefore
\begin{equation}\label{eq:largec}
 r\ge2,\qquad k\ge8XY^2+4,\qquad c>kY>8XY^3>2R.
\end{equation}
Since $P_0>a_0$ and $c>k$, monotonicity also gives $p>r$ and $c>2p$.

The strong norm yields $f=\chi_{a_0}(m)$ and $I_*=\psi_{a_0}(m)$.
From $c^2\mid I_*$, strong divisibility gives $p\mid m$. Writing $m=pu$
and using composition gives
\[
 I_*=c\psi_D(u),\qquad D^2\equiv1\pmod c.
\]
The recurrence implies $\psi_D(u)\equiv uD^{u-1}\pmod c$.
Because $D$ is a unit modulo $c$, one gets $c\mid u$, so
\begin{equation}\label{eq:rank}
 pc\mid m,\qquad c\mid m,\qquad m>2p.
\end{equation}

The auxiliary equation gives an index $\nu\ge1$ with
$S|V|=\chi_S(\nu)$ and $y=\psi_S(\nu)$. It is odd: the even-index
abscissas are $\pm1$ modulo $S$, whereas $S|V|$ is divisible by $S$.
Write $\nu=2u+1$. The integer polynomials
\[
 L_0(Z)=1,\quad L_1(Z)=4Z-3,\quad
 L_{u+1}(Z)=(4Z-2)L_u(Z)-L_{u-1}(Z)
\]
satisfy, by induction,
\begin{align*}
 \chi_b(2u+1)&=bL_u(b^2),&L_u(0)&=(-1)^u(2u+1),\\
 L_u(1-a_0^2)&=(-1)^u\psi_{a_0}(2u+1).
\end{align*}
Here $Q\equiv1-a_0^2\pmod f$, $V\equiv-c\pmod f$,
$Q\equiv0\pmod c$ and $V\equiv-R\pmod c$.
It follows that $\psi_{a_0}(\nu)\equiv\pm c\pmod f$.
Squaring and using $\chi_{a_0}(2j)=1+2\Delta\psi_{a_0}(j)^2$ gives
\[
 \chi_{a_0}(2\nu)\equiv\chi_{a_0}(2p)\pmod{\chi_{a_0}(m)}.
\]

The signed step-down lemma says that, for $0<j<m$, a congruence
$\chi_{a_0}(t)\equiv\chi_{a_0}(j)\pmod{\chi_{a_0}(m)}$
forces $t\equiv\pm j\pmod{2m}$. To see this, use the addition
identities to reduce and reflect the index to a signed
$\chi_{a_0}(s)$, $0\le s\le m$. The endpoint $s=m$ has residue zero.
A negative representative cannot equal the specified positive one,
since $2\chi_{a_0}(m-1)<\chi_{a_0}(m)$; among positive
representatives strict monotonicity forces $s=j$.
Apply it with $j=2p<m$. Then $\nu\equiv\pm p\pmod m$, hence
modulo $c$. Reducing $L_u(Q)$ modulo $c$ separately gives
$\nu\equiv\pm R\pmod c$. Thus $p\equiv\pm R\pmod c$.
The strict inequalities $0<p,R<c/2$ force
\begin{equation}\label{eq:pR}
 p=R.
\end{equation}
This is the point at which the old main index is recovered; the omitted
quotient has not appeared in any premise.

\subsection{How the residue carrier restores the quotient}

For $P_{13,70}$ and $P_{13,71}$, let $M=E$; for $P_{13,72}$, let
$M=q^4$. In each case $t=Mn-k+R$, $P_0\equiv1\pmod M$, and
$M\mid S^2$. The odd-ordinate recurrence gives
\begin{equation}\label{eq:oddres}
 \psi_S(2u+1)\equiv(-1)^u\pmod{4S^2}.
\end{equation}
Also $\tau\equiv1\pmod M$ and $k\equiv2r\pmod M$.
Consequently $t=\pm\tau y\equiv\pm1\pmod M$, and for a sign
$\varepsilon\in\{1,-1\}$,
\[
 2r-R\equiv\varepsilon\pmod M.
\]
Equations~\eqref{eq:largec}--\eqref{eq:pR} give $1\le r<R<M$.
Thus $|2r-R-\varepsilon|<R<M$, forcing the exact equality
$2r=R+\varepsilon$.

The negative sign is impossible. If $R=2r+1$, put $B_0=2a_0^2-1>P_0$.
Composition and monotonicity give
\[
 \psi_{a_0}(2r)=2a_0\psi_{B_0}(r)
     \ge2a_0\psi_{P_0}(r)=a_0k>k(Y+1).
\]
Then $c=\psi_{a_0}(R)>k(Y+1)$, contradicting
$c=kY+\eta<k(Y+1)$. Therefore $2r=R+1$.

The full first-Pell recurrence supplies $k\equiv2r\pmod{4EY}$,
even in the smaller-modulus variant. In every case we can restore
\begin{equation}\label{eq:hrestore}
 h=\frac{k-R-1}{E}\in\mathbb Z_{>0}.
\end{equation}
Positivity follows from $k\ge8XY^2+4>R+1$.
Together with $|t|$ and the appropriate positive rescaled multiplier,
these are a zero of the exact retained parent. Its native/refutation
calibration gives a finite syntactic ZFC refutation.

Conversely, a refutation supplies canonical outer parent witnesses with
$R\equiv3\pmod4$, $A=\tau^2$, $c=\psi_{a_0}(R)$ and $N_k=N_t=N_1=1$.
The finite matrix-period argument permits the strong multiplier to be
chosen divisible by $\tau E$, $AE$, or $Aq^2$, respectively. Choose
fresh auxiliary witnesses for this new strong pair:
\[
 V_*=\frac{\chi_S(R)}{S},\qquad y_*=\psi_S(R),\qquad
 T_*=\frac{V_*+c+Rf^2}{cf^2}.
\]
The odd quotient identities give $V_*\equiv-c\pmod{f^2}$ and
$V_*\equiv-R\pmod c$. Since $\gcd(f,c)=1$ and $f^2\equiv1\pmod c$,
$T_*$ is a positive integer.
Equation~\eqref{eq:oddres} gives $y_*\equiv-1\pmod M$.
Thus $t_*:=\tau y_*$ is $-1$ modulo $M$, while $k-R$ is 1 modulo $M$.
The number $n=(t_*+k-R)/M$ is a positive integer, and every factor
in the new equation is $+1$. This proves the converse existence map.

\subsection{The rational carrier and twelve positive unknowns}

The following inherited descent is included because it is the exact
integrality mechanism required by \eqref{eq:p12}. It is substantially
stronger than merely clearing a denominator in an arbitrary Pell equation.

\begin{lemma}[Adjacent-coefficient integrality]\label{lem:integral}
Let $Q\ge2$ and $V$ be integers. If $y\in\mathbb Q$ satisfies
$QV^2-(Q-1)y^2=1$, then $y\in\mathbb Z$.
\end{lemma}
\begin{proof}
Replace $V,y$ by their absolute values and write $y=b/d>0$ in lowest
terms. Neither coordinate is zero, and $d^2\mid Q-1$. If $V=1$, then
$y=1$. Otherwise $y>V>1$. Put $D_0=Q-1$ and
\[
 V'=(2Q-1)V-2D_0y,\qquad
 y'=\bigl|(2Q-1)y-2QV\bigr|.
\]
The first coordinate is integral. The identity
\[
 ((2Q-1)V)^2-(2D_0y)^2=V^2+4D_0>0
\]
shows $V'>0$, while $y>V$ gives $V'<V$.
Direct expansion preserves the same norm equation. The reduced
denominator remains exactly $d$, since $2Q-1\equiv1\pmod d$.
Descent in the positive integer $V$ reaches $V=1$, where the denominator
is one. Thus the original $d$ was one.
\end{proof}

At a zero of \eqref{eq:p12}, the unit arguments give $N_1=N_t=1$
and $B_{12}=0$. Therefore \eqref{eq:Rbounds} holds. Since
$e=k-R-1\ge1-R>-E$ and $n\ge1$, we have $Z>0$.
The positive rational $t=Z/E=n+e/E$ satisfies
\begin{equation}\label{eq:ratfold}
 (Q-1)t^2=AF_2(QV^2-1).
\end{equation}
Here $A\mid Q$ and $F_2\equiv1\pmod A$. Also
\[
 Q-1=\Delta F_2-a_0^2,
 \qquad QV^2-1\equiv-D^2\pmod{F_2}.
\]
The anchors in $m=a_0Di$ ensure $F_2\equiv1\pmod{a_0}$ and
$F_2\equiv1\pmod D$. Consequently
\[
 \gcd(A,F_2)=1,\qquad
 \gcd\bigl(AF_2,(Q-1)(QV^2-1)\bigr)=1.
\]
Prime valuations in \eqref{eq:ratfold} show that both $A$ and $F_2$
have even valuations. This use of valuations is valid for rational
$t$, whose prime valuations are still integers. Hence
$A=\tau^2$ and $F_2=f^2$ for positive integers $\tau,f$.

Now $Q=(\Delta\tau mc^2)^2>1$, and
$y=t/(\tau f)>0$ satisfies the hypothesis of
Lemma~\ref{lem:integral}. It follows that $y$, and hence $t$, is an
integer. Thus $E\mid e$. Since $e>-E$, one has $e\ge0$.
The possibility $e=0$ would give $k=R+1$. The first Pell equation
forces either $k=2$ or $k\ge8XY^2+4$, both incompatible with
\eqref{eq:Rbounds}. Therefore $h=e/E>0$. These restored $h,t$ give
an exact zero of the anchored thirteen-variable parent, proving soundness.

For completeness, retain the canonical outer data of a refutation,
so $e=hE$ with $h>0$. The parent's matrix-period construction supplies
arbitrarily large strong multipliers divisible by $\tau a_0D$.
Recompute the auxiliary pair after this choice. Its combined root
$t=\tau f y$ can be made larger than $h$. Then $n=t-h>0$, $Z=Et$,
and $B_{12}$ is exactly $E^2$ times the vanishing parent residual.
Thus \eqref{eq:p12} has a positive zero.

\subsection{Exact degrees, operation counts, and logical conclusion}

Deleting the old $N_k$ removes its three operations and one multiplication
in the final product. The first two residue carriers add one multiplier
scaling and three operations for $t$, preserving the parent operation
counts. The third uses a new $q^4$ gate and a lower-degree multiplier.
The twelve-variable formula clears the single rational carrier by $E^2$;
its literal graph has 74 gates. Counts concern the supplied complete
programs, with fixed integer leaves and reused values free.

Exact degree verification accounts for the cancellation
\[
 N_1=2ac(X+\gamma H)+(X+\gamma H)^2-Hc^2,
\]
which makes $\deg N_1=18$. Merely propagating syntactic degrees
through $D^2-\Delta c^2$ would overestimate it. Each output has a
nonzero coefficient at the claimed top degree:
\begin{center}\small
\begin{tabular}{@{}lrl@{}}
\toprule
Equation & Degree & One attained coefficient\\
\midrule
$P_{13,70}$ &206&$-24\beta^{124}$\\
$P_{13,71}$ &162&$-24\beta^{96}$\\
$P_{13,72}$ &154&$-24\beta^{92}$\\
$P_{12,74}$ &438&$+24\beta^{272}$\\
\bottomrule
\end{tabular}
\end{center}
The full monomials, homogeneous leading forms, exact graph checks and
coefficient bindings are included in the verification artifacts.

For each of these four fixed polynomials $P$, the existence maps and
the retained native/refutation calibration prove, in ordinary arithmetic
formalizable in ZFC,
\[
 \mathrm{ZFC}\vdash
 \bigl(\neg\exists\mathbf x\in\mathbb Z_{>0}^{\,d}\;P(\mathbf x)=0\bigr)
 \ \longleftrightarrow\ \Con(\mathrm{ZFC}).
\]
The positive domain is part of the statement. It is not a claim about
unrestricted integer zeros of the same unmodified expression.

% END compact_proof.tex


\section{A direct quartic from pooled powers and integral projection}\label{sec:quartic}
% BEGIN projection_proof.tex
The direct branch keeps every quantified coordinate unrestricted over
$\mathbb Z$ and uses only small integer coefficients. Its principal
equation is
\begin{equation}\label{eq:principal}
 \boxed{\displaystyle
 P_{\DirectVars}(x_0,\ldots,x_{\DirectLast})
   =\sum_{j=0}^{\DirectRowLast}r_j(\mathbf x)^2=0,
 \qquad\mathbf x\in\mathbb Z^{\DirectVars}.}
\end{equation}
Every residual is the explicit integer quadratic printed in
Appendix~\ref{app:rows}. The archive also supplies every coefficient of
the fully collected quartic. A JSON term $[c,[i_1,\ldots,i_k]]$ means
$cx_{i_1}\cdots x_{i_k}$; an empty index list means a constant.
The list contains all quantified coordinates. There are no unnamed
existential witnesses.

This construction combines three separately attributed ingredients:
the factorial-free arithmetic source with the affine radix described
below, the strict local ranges and fixed-base pooling of the incoming
Pooled Powers manuscript \cite{pooled}, and the integral quadratic
projection of the incoming 234-variable manuscript \cite{projection}.
The new deliverable is their compatible, fully emitted composition.
The earlier methods and the already submitted diagonal growth choice
are not presented as discoveries of this report.

\subsection{Two quadratic rows certify a block of powers}

The inherited signed Pell frame provides integers $W\ge258$ and $A$
with $|A|>W^W$. For a block of $r\le15$ requests with the same
nonzero integer base $n$, it supplies canonical signed Pell pairs
$(X_i,Y_i)$ at signed indices $k_i$ and coordinates $Z_i=AY_i$.
Every base, proposed output and index has square less than $W$.
Put
\[
 T=2An-n^2-1,\qquad D_i=m_i-X_i+Z_i,
\]
and use the following quantities only in the proof:
\[
 H=\sum_iY_i(m_i-X_i),\qquad
 S=\sum_i(D_i^2-Y_i^2).
\]
Instead of a separate quotient for every requested power, introduce
just two integer coordinates $K,L$, with the quadratic rows
\begin{align}
 K-2\sum_iY_i(m_i-X_i)-(2A-n)L&=0,\label{eq:pool1}\\
 \sum_i\bigl((m_i-X_i+Z_i)^2-Y_i^2\bigr)-nK+L&=0.
 \label{eq:pool2}
\end{align}
No coordinate for $H$, $S$, or $T$ is needed.

Eliminating $K$ in the proof gives $S-2nH=TL$. The exact identity
\[
 \sum_i(D_i-nY_i)^2=S-2nH-T\sum_iY_i^2
\]
therefore makes the left side divisible by $T$.
Let $b=\lfloor\sqrt{W-1}\rfloor$. Then
$|n|,|m_i|,|k_i|\le b$, and the Pell congruence, multiplied by
$n^b$ to clear every negative index, makes $T$ divide
\[
 E_*:=\sum_i\bigl(m_i n^b-n^{b+k_i}\bigr)^2.
\]
All exponents here are nonnegative proof notation. The estimates
\[
 0\le E_*\le4r b^{4b}<W^{2b+1}\le W^W<|T|
\]
force $E_*=0$. Each square is zero. If $|n|\ge2$, a negative
$k_i$ would give the impossible integer identity $m_i n^{-k_i}=1$.
Thus all $k_i\ge0$ and $m_i=n^{k_i}$.

Conversely, genuine powers and canonical Pell pairs make each
$D_i-nY_i$ divisible by $T$. The identity makes
$L=(S-2nH)/T$ integral, and \eqref{eq:pool1} defines integral $K$.
This proves both directions of the pooled certificate. Its proof uses
the existing frame; no larger frame or new variable-exponent operation
has been inserted into the polynomial.

\subsection{Strict local ranges and the sign order}

The incoming strict-horizon argument chooses a cofinal proof horizon
strictly larger than all finitely many local labels and parameters.
Consequently its local binomial range tests may enforce $0\le u<t$
instead of $0\le u\le t$. In the literal source, the local range
denominator is changed from $q^{t+1}$ to the already required $q^t$.
The four now private coordinates and their rows can be removed with
the stated restoration identities. This is an existential completeness
change justified by choosing a larger horizon; it is not an equality
of arbitrary old and new coordinate tuples.

After this change the source has eleven requests at $q+1$ and three at
$q$. The three other requests remain individual. Applying
\eqref{eq:pool1}--\eqref{eq:pool2} to the two blocks gives a literal
223-coordinate system of 160 quadratic residuals. The finite certificate
authenticates the exact affine offsets of all 14 requests, all retained
$AY$ and Pell norm rows, and every base/index/output in the square frame.

In our arithmetic source $q=6z+3$ holds before any sign inference.
Hence $|q|\ge3$ and $|q+1|\ge2$ for every integer $z$. The block
theorem immediately restores all 14 individual powers and their
nonnegative indices, after the inherited orientation-first theorem
has supplied the canonical pairs. Thus the sign bootstrap in the next
sections remains acyclic. Completeness starts from genuine powers and
uses the converse part of the block theorem to construct the four new
pooled coordinates.

\subsection{Integral projection along a common quadratic direction}

Here is the exact algebra behind the final six projections. It also
applies to the separately retained unchanged 225-variable source.

\begin{lemma}[Quadratic projection with a unit slope]\label{lem:projection}
Let $Q_0,\ldots,Q_{r-1}\in\mathbb Z[x_0,\ldots,x_{n-1}]$ have
degree at most two. Suppose $v\in\mathbb Z^n$ is in the common kernel
of their quadratic Hessians, $v_p=1$, and
$\partial_vQ_j=a_j\in\mathbb Z$ for every $j$.
If $a_s=\pm1$, the common integer zero set is in polynomial bijection
with the common zero set in $n-1$ variables of
\[
 R_j(y)=\left.(Q_j-a_sa_jQ_s)\right|_{y_p=0},\qquad j\ne s.
\]
Every $R_j$ still has degree at most two.
\end{lemma}
\begin{proof}
The common-kernel condition gives $Q_j(y+vt)=Q_j(y)+a_jt$.
For $y_p=0$, set $T(y)=-a_sQ_s(y)$ and $x=y+vT(y)$.
All $Q_j(x)$ vanish whenever all $R_j(y)$ vanish.
Conversely, for a common integer zero $x$ put $y=x-vx_p$.
Then $y_p=0$, and $Q_s(x)=0$ forces $x_p=T(y)$.
The two polynomial maps are inverse on the zero sets. No division,
positivity condition, or unproved divisibility is involved.
\end{proof}

The six literal certificates apply Lemma~\ref{lem:projection} to the
223-coordinate source and give \DirectVars\ coordinates and
\DirectRows\ residuals. Their independently written checker obtains
each reduced row directly by the displayed subtraction identity,
rather than invoking the producer's change-of-coordinates code.
It checks the full direction, every slope, the unit condition, every
new coefficient, and the inverse maps.

Finally elementary integer row operations $Q_i\leftarrow Q_i+kQ_j$
reduce coefficient height and select useful support/height tradeoffs.
The inverse uses $-k$ in reverse order. Such changes preserve
simultaneous vanishing, while changing the sum-of-squares polynomial
itself. The selected polynomial is consequently identified by its own
full coefficient file and cryptographic hash.

Because a sum of integer squares vanishes exactly when every residual
vanishes, these explicit zero-set maps and the arithmetic/refutation
calibration prove the equivalence claimed for \eqref{eq:principal}.

% END projection_proof.tex


% BEGIN diagonal_background.tex
\section{One diagonal power serves three purposes}\label{sec:diagonal}

The diagonal choice $F=t^t$, the alias $R=F$, and the linear odd beta
spacing already occur in the incoming Linear CRT manuscript~\cite{linearcrt}.
We give the needed proofs here to make the new sign-bootstrap and pooling
composition reviewable. The distinct feature of this source is the fixed
radix class $q\equiv3\pmod6$.

The direct arithmetic construction encodes a finite multiset of valid
instances of a single-conclusion proof calculus. An exact balance between
conclusion labels and premise labels characterizes a proof of the target.
The existing compiler represents a finite sequence of multiplicative
accumulator states through beta remainders and certifies its local rows
through amplified binomial CRT conditions. We retain that architecture
and change the growth parameter and sign bootstrap.

Throughout this section, $t\ge2$ is an even integer. In the actual
polynomial source, $t=16+2h^2$, so these premises are automatic. Put
\[
 F=t^t,\qquad L_t=\lcm(1,2,\ldots,t).
\]
The variable exponent is proof notation: $F=t^t$ is enforced by an
existing simultaneous Pell power certificate in the ordinary source.
It is not an operation in the final quartic.

\subsection{Separating bounded multisets at one integer}

For a finite list $A=(a_1,\ldots,a_m)$ with $1\le a_i\le t$ and
$0\le m\le t$, define the auxiliary polynomial
\[
 H_A(X)=\prod_{i=1}^{m}(1+a_iX).
\]
The empty product is one. This polynomial is used in the proof only;
the length of the list is not the number of final unknowns.

\begin{lemma}[Diagonal evaluation separates the lists]\label{lem:labels}
If $A$ is such a list, $B$ is any finite list of labels in
$\{1,\ldots,t\}$, and $H_A(F)=H_B(F)$ with $F=t^t$, then $A$ and $B$
have the same multiplicities of every label. Only one initial length
bound is required.
\end{lemma}
\begin{proof}
First their lengths are equal. If $m>n$, then $m\ge n+1$ and
\[
 H_A(F)\ge(F+1)^m\ge(F+1)^{n+1}
       >(t(F+1))^n\ge(tF+1)^n\ge H_B(F).
\]
The strict inequality follows from $F+1>t^n$. The following weak
inequality also covers $n=0$. If $n>m$, the symmetric comparison uses
$m\le t$ and $F+1>t^m$; no prior upper bound on $n$ is needed.

Let the common length be $m$. For $0\le k\le m$, the coefficient of
$X^k$ in either product is a positive integer and is at most
\[
 \binom{m}{k}t^k\le\binom{t}{k}t^k
 \le t^{t-k}t^k=t^t=F.
\]
Here $\binom{t}{k}=\binom{t}{t-k}\le t^{t-k}$.
Thus every coefficient of $H_A-H_B$ has absolute value strictly less
than $F$. If a nonzero coefficient had least index $k$, division of
$(H_A-H_B)(F)=0$ by $F^k$ and reduction modulo $F$ would force that
coefficient to be divisible by $F$, a contradiction. Therefore
$H_A=H_B$ in $\Z[X]$.

Each $1+jX$ is a primitive linear polynomial, and two such factors with
positive $j$ are associates only when their labels agree. Unique
factorization, or comparison of the rational roots $-1/j$, now gives
the asserted multiplicities.
\end{proof}

Truth label zero has weight one and is deliberately omitted from these
lists. The retained proof-flow interface treats truth as freely available.
The lemma therefore provides exactly the projected label balance needed
by that interface. It does not claim that all integers $1+jF$ are pairwise
coprime; they need not be.

\subsection{The same power covers the CRT rows}

\begin{lemma}[Elementary least-common-multiple bound]\label{lem:lcm}
For even $t\ge2$, $L_t^2\le t^t=F$.
\end{lemma}
\begin{proof}
Write $L_t$ as the product of the largest power of each prime not
exceeding $t$. Each factor is at most $t$. There are at most $t/2$
such primes: besides 2, they belong to the $t/2-1$ odd integers from
3 through $t-1$. Thus $L_t\le t^{t/2}$, and squaring proves the claim.
\end{proof}

The inherited factorial-free CRT theorem uses
\[
 R\ge L_t^2,\qquad B=\binom{R+t}{t}.
\]
It proves two different prime-power facts. Every row $i\in\{1,\ldots,t\}$
has a full prime power $p^e\Vert B$ with $p^e\mid R+i$ and $p^e>t$;
conversely every full prime power of $B$ divides at least one such row.
The first fact supplies row-by-row soundness. The second supplies a
consistent reverse CRT construction at all prime factors of $B$.
The retained proof includes their valuation arguments and the amplified
range extraction; this report does not replace them by an assumption of
pairwise coprimality of the consecutive factors.

Lemma~\ref{lem:lcm} permits the literal specialization
\begin{equation}\label{eq:shift}
 R=F=t^t,\qquad B=\binom{F+t}{t}.
\end{equation}
The former square relation $R=F^2$ is unnecessary. This is a change to
the semantically certified source, followed by a polynomial alias and
exact elimination of its now redundant coordinate. It is not the
unjustified deletion of an equation from the old zero set.

\subsection{The required parity is already present}

The same choice also provides
\[
 v_2(F)=t\,v_2(t)\ge t>v_2(t!).
\]
Consequently an arbitrarily large odd multiple of $F$ can be chosen
divisible by $t!$: multiply by an odd multiple of the odd part of $t!$.
This is precisely the freedom used to choose a beta spacing large enough
for every pending proof-accumulator state. The polynomial need not
represent $t!$; it occurs only in the constructive existence proof of an
appropriate spacing.

Also $B=\binom{F+t}{t}$ is odd. For $1\le i\le t$,
$v_2(i)<v_2(F)$, so $v_2(F+i)=v_2(i)$. The powers of two in the
numerator product and in $t!$ therefore cancel exactly. This parity
allows the subsequent CRT representatives and beta codes to be chosen in
the required residue classes.

\section{Recovering signs before arithmetic interpretation}\label{sec:signs}

The direct final equation has unrestricted integer unknowns. Positivity
must therefore follow from equations, and the order in which it follows
matters. We replace the old positive-square radix by
\begin{equation}\label{eq:radix}
 q=6z+3,\qquad H=3+6\sum_{j=1}^{9}a_j,\qquad q=(BH)^x,
\end{equation}
where the $a_j$ are the nine actual CRT representatives. The last identity
is an existing polynomially certified power request.

Before the sign of any representative or of $B$ is known,
\[
 |q|\ge3,\qquad |q+1|\ge2.
\]
Thus the signed power theorem applies to both bases. The requests with
base $q+1$ force the nine representative indices to be nonnegative;
the request with base $q$ and index $y+1$ gives $y+1\ge0$.
It follows that $H\ge3$ and $y\ge-1$ without assuming a binomial value
or a positive radix.

The beta spacing can now be linear in its cofinal parameter:
\begin{equation}\label{eq:beta}
 y=F(2v+1).
\end{equation}
Because $F=t^t\ge2$ and $y\ge-1$, its odd multiplier cannot be negative.
Hence $y\ge F$ and $v\ge0$. This recovers the sign that a square formerly
imposed, while retaining a cofinal collection of legal spacings. The full
bootstrap, including the normalization ratios, the signed first binomial
digit, $B>0$, $x>0$, and finally $q>0$, is given in the following proof.

% BEGIN direct_proof.tex
\subsection{The signed power and remainder interfaces}

The unchanged semantic anchor and full orientation frame supply
$W\ge258$, $|A|>W^W$, canonical signed Pell pairs, and bounds on the
actual power bases, indices and outputs. The two pooled blocks are
covered by the theorem just proved. For the three individual requests,
the usual congruence
\[
 X_k(A)+(n-A)Y_k(A)\equiv m\pmod{2An-n^2-1}
\]
gives the same conclusion when $|n|\ge2$: the bounded congruence
is the equality $m=n^k$, and a negative $k=-j$ would force
$mn^j=1$. This is impossible for an integer $m$. The bounds use
absolute values and do not require $n$ to be positive.

The common normalized remainder certificate is
\[
 M^2=1+4\sum_j R_j^2+a^2+b^2+c(c+1).
\]
As $c(c+1)\ge0$ for every integer $c$, one has $M\ne0$ and
$|R_j|<|M|/2$. If $L_jv_j=M$ and $v_jU_j=MQ_j+R_j$, then
$L_j,v_j\ne0$, $v_j\mid R_j$, and the integer $r_j=R_j/v_j$ satisfies
\[
 U_j=L_jQ_j+r_j,\qquad |r_j|<|L_j|/2.
\]
A congruent remainder in this open centered interval is unique, even
for negative $L_j$. The first binomial interpretation can therefore be
recovered before the radix has been proved positive.

\subsection{An acyclic bootstrap}

The horizon $t=16+2h_0^2$ is positive and even. Its individual power
request first gives $F=t^t\ge2^{64}$.
The fixed class $q=6z+3$ makes the two pooled bases nonzero nonunits.
Their power requests give $x\ge0$, $y+1\ge0$, and nonnegative values
of all nine actual representatives $a_j$. Hence $H\ge3$.
From $y=F(2v+1)$ and $y\ge-1$ follows $y\ge F$.

The raw index equation is
\[
 B\,u=C_0+F+1,
\]
where $C_0$ is the already nonnegative actual index representative.
Thus $B\ne0$ and $|BH|\ge3$. The individual self-radix request
gives $q=(BH)^x$. Since $q\ne1$, one has $x\ge1$.
The first and last beta endpoint equations are
\[
 x=F+1+(y+1)u_0,
 \qquad x=1+\bigl(1+(t+1)y\bigr)u_t.
\]
If $u_0\le-1$, the first gives $x\le F-y\le0$, impossible.
Therefore $x\ge F+1$. The last endpoint then forces $u_t\ge1$, so
\begin{equation}\label{eq:xlarge}
 x\ge2+(t+1)y\ge2+(t+1)F.
\end{equation}
Write $Q=|q|$ and $b=|B|$. Without assuming either sign, we now have
\[
 Q=(bH)^x\ge2^x,\qquad Q\ge9b^2,\qquad b\le Q/9.
\]
All low-digit denominator exponents are nonnegative and at most $x$.
Their exact ratios to the common modulus $M=q^x$ are integral.

\subsection{The first binomial digit forces a positive radix}

For $0\le j\le t$, put $C_j=\binom{F+t}{j}$.
Since $F\ge t$, these coefficients increase up to $t$, and
\[
 C_j\le C_t=\prod_{i=1}^t\left(1+\frac{F}{i}\right)
       \le(F+1)^t=:V_*.
\]
The bound $F+1\le2^{F-1}$ and \eqref{eq:xlarge} imply
\[
 \frac{V_*}{Q}\le2^{-F-t-2},\qquad \frac{V_*}{Q-1}<1/32.
\]
The authenticated numerator is $U=(q+1)^{F+t}$, and its digit row is
\[
 U=q^t(B+q\xi)+r,\qquad |r|<Q^t/2.
\]
For either sign of $q$, the binomial theorem gives the canonical
integer remainder
\[
 r_0=\sum_{j=0}^{t-1}C_jq^j,
 \qquad |r_0|<\frac{V_*Q^t}{Q-1}<Q^t/32.
\]
Centered remainder uniqueness yields $r=r_0$, and division by $q^t$
gives $B\equiv C_t\pmod q$. But $|B|\le Q/9$ and
$C_t<Q/32$, so $|B-C_t|<Q$. Thus
\[
 B=\binom{F+t}{t}>0,\qquad q=(BH)^x>0,\qquad N=B^x>0.
\]
Because $t\ge2$, this binomial coefficient exceeds $F+1$.
The earlier divisibility $B\mid C_0+F+1$ excludes $C_0=0$;
hence $C_0\ge1$, as required by the row-index extraction.
This argument handles the entire negative-radix possibility before
excluding it.

\subsection{The two prime-power covering statements}

For completeness, let $\Lambda=\lcm(1,\ldots,t)$, $R\ge\Lambda^2$,
and $B=\binom{R+t}{t}$. Every row $i\in\{1,\ldots,t\}$ has a full
prime power $p^e\Vert B$ dividing $R+i$ with $p^e>t$.
Indeed, if every $v_p(R+i)\le2\lfloor\log_p t\rfloor$, then
$R+i\mid\Lambda^2$, contradicting $R+i>\Lambda^2$.
Choose $p$ with $v=v_p(R+i)>2a$, $a=\lfloor\log_p t\rfloor$.
For $j\ne i$, the distance $|j-i|\le t$ implies
$v_p(R+j)=v_p(j-i)$. Therefore
\[
 e=v_p(B)=v-v_p\!\left(t\binom{t-1}{i-1}\right).
\]
In the Legendre sum for the subtracted valuation, each level up to $a$
contributes zero or one, and all later levels contribute zero.
Thus $v-a\le e\le v$, which proves $p^e>t$ and $p^e\mid R+i$.

Conversely, every full prime power of $B$ divides at least one row.
For a fixed prime let $v=\max_i v_p(R+i)$. Legendre's formula for
$v_p(B)$ is a sum of zero-or-one interval-carry terms, with no terms
above level $v$. Hence $v_p(B)\le v$. Both statements are needed:
one supplies a large row selector, and the other covers every prime
factor during reverse CRT construction. Consecutive $R+i$ need not
be pairwise coprime.

Here $R=F=t^t\ge\Lambda^2$ by Lemma~\ref{lem:lcm}, so the
retained factorial-free binomial CRT compiler applies.

\subsection{Remaining local bounds and the beta endpoint}

The inherited local-polynomial ledger uses exponent $2^{23}$.
For $R_{\rm maj}=2+t+x+y+F$, the established bounds give
$R_{\rm maj}<2x$ and $x\ge2^{64}+1$. The retained uniform estimate
$(2x)^{2^{23}}<2^x$ therefore applies. Its derived-label degree
bound $2^{22}(t+1)<x$ also holds. Since $H=3+6\sum a_j$ dominates
the old representative majorant, the other ten binomial digit
remainders retain their strict relative margin below $1/32$.
Their exact divisibilities give the same local polynomial constraints,
with the strengthened strict local ranges from the pooling source.

For each row the selected prime power $p^e>t$ identifies the correct
index. The amplified range lemma, the precision $x$, and the bound
$0\le L<2^x\le p^{ex}$ turn the local congruence into $L=0$.
Thus all supplied Horn-rule and accumulator rows are genuine.

There is a small boundary case to settle. If $y=F$, the first beta
endpoint gives initial proper remainder zero, while the last gives
final remainder one. The recovered local equations are
\[
 S_{i+1}W_{\mathrm{concl}(i)}
  =S_iW_{\mathrm{prem}_1(i)}W_{\mathrm{prem}_2(i)},
 \qquad W_j=1+jF>0.
\]
A zero initial state propagates to every state, contradicting the last.
Hence $y>F$. The actual first remainder is $F+1$, the last is one,
and all intervening states are positive. Their product identity and
Lemma~\ref{lem:labels} recover the exact multiplicities of proof labels.
The retained single-conclusion flow criterion then yields a $C_{222}$
refutation. It allows cycles and repeated premises: summing the balance
over labels outside the finite closure of truth gives the usual
contradiction if the target is outside that closure.

\subsection{Constructive completeness}

From a finite $C_{222}$ refutation, choose a finite proof tree and an
even horizon $t=16+2h_0^2$ strictly larger than its size, labels and
every local parameter. Pad with truth to $t$ rows. Set $F=t^t$ and
form the finite pending-multiset accumulator path with weights $1+jF$.

Let $O$ be the odd part of $t!$. Choose an arbitrarily large odd
$u=O(2m+1)$, and set $y=Fu$, $v=(u-1)/2$. The inequality
$v_2(F)\ge t>v_2(t!)$ gives $t!\mid y$. Choose $u$ large enough
that $y>F$ and exceeds every accumulator value. The beta moduli
$1+(i+1)y$ are pairwise coprime and odd. CRT gives a positive beta
code $x$, cofinal in its residue class.

Set $B=\binom{F+t}{t}$ and $N=B^x$. The two prime-power statements
assign compatible genuine rows at every prime factor of $B$.
The inherited precision-$N^2$ CRT construction supplies all nine
nonnegative representatives and range witnesses. Since $B$ is odd,
$N^2$ is odd, so their cofinal lifts also permit the parity and
dominance choices for the semantic anchor.

Now $H=3+6\sum a_j$ is positive, odd and divisible by three.
Thus $q=(BH)^x$ is positive, odd and divisible by three, and
$z=(q-3)/6$ is an integer. The standard canonical construction gives
the individual powers, Pell frame, normalized remainders and common
norm completion. The converse pooling theorem supplies its four new
coordinates. All source rows vanish. The exact projection and row
maps then give a zero of the final quartic.

Together with soundness and the inherited logical calibration, this
proves that the new arithmetic source, its pooled form, and its final
integer quartics have zeros exactly when there is a finite syntactic
ZFC refutation.

% END direct_proof.tex



% END diagonal_background.tex


\section{Exact finite verification}\label{sec:verification}
% BEGIN finite_verification.tex
The release separates three kinds of evidence. The ordinary proofs give
the unbounded implications between zeros and refutations. Exact finite
checks bind those proofs to literal source expressions and inverse maps.
Cryptographic hashes identify the submitted files and their inherited
sources. None of these is represented as a substitute for another.

For the compact equations, a checker independent of their producer
verifies the whole directed acyclic program, including the declared
positive coordinate domain, topological order, and occurrence of every
coordinate, constant and operation in the output. Exact sparse-polynomial
identities at intermediate cuts cover every gate. Full homogeneous
leading forms certify the degrees, accounting for cancellation in the
main Pell norm. Specified nonzero top coefficients certify attainment
at the actual fixed tuple. The four complete integers are independently
reconstructed from the finite indexed-tile recipe and compared byte for
byte with their canonical payloads. Parent substitutions and deliberate
arithmetic and metadata corruptions are also checked.

For the direct source, the earlier source-level checker authenticates
the changed horizon, growth, radix, beta and binomial-shift rows; every
retained power request and frame atom; Horn factorizations; monic integer
projections; coordinate translations; and the recorded affine maps.
The separate strict/pool checker imports no producer and checks the
literal strict-range edits, all fourteen pooled power identities, their
Pell and $AY$ rows, and the exact 36-square frame containing all seventeen
remaining power requests. It expands the new pooled residuals exactly.

The integral projection checker reconstructs each reduced row by the
subtraction formula of Lemma~\ref{lem:projection}. It verifies every
common quadratic direction, unit slope, inverse restoration and new
coefficient. Reversible row operations are replayed forwards and
backwards. Independent squaring and collection reproduces every final
quartic coefficient, its degree, support and height. The 220-variable
companion starts from the unchanged incoming Pooled Powers source,
whose own complete release verification and deterministic rebuild
were also replayed successfully.

Normal and optimized Python runs are included; the check conditions use
explicit exceptions, so optimization does not silently remove the
verification gates. Fresh runs from a different working directory test
portability. A release-level checker verifies the selected equation
copies, printed residual data and file manifest and invokes the component
checks. The standard library suffices for the mathematical checks.

The checks are intentionally finite. They do not materialize an alleged
refutation, infer correctness from testing bounded candidate zeros,
or emit a new end-to-end Lean or Rocq theorem. Sample evaluations and
negative controls support implementation auditing; the symbolic
identities and ordinary proofs carry the mathematical claims.

% END finite_verification.tex


\section{Calibration, reproducibility, and limits}\label{sec:status}

The direct arithmetic encodes the numerical calculus $C_{222}$ itself.
Its finite constructor polynomials and full rule list are retained with
the original proof sources. Its logical interfaces are
\begin{align*}
 C_{222}(1)&\longleftrightarrow\text{a finite ZF refutation},\\
 \text{a finite ZF refutation}&\longleftrightarrow
                  \text{a finite ZFC refutation}.
\end{align*}
The repository declarations are
\file{KEndpoint.c222_iff_zf_refutation} and
\file{ConstructibleChoice.zf_refutation_iff_zfc_refutation}.
The latter is the constructible-universe equiconsistency interface.
No K215 simulation or Jones universal polynomial is used on this route.

The compact branch retains the entire native indexed943/radix67 tuple
and its calibrated finite-refutation interface. This branch does retain
the inherited K215/native encoding. Constants from different native
tuples are not interchangeable. Its four supplied values are fixed
explicit integers, rather than parameters awaiting a consistency-dependent
choice or unspecified coefficients chosen by an existence theorem.

The new reasoning is an ordinary mathematical proof formalizable in ZFC.
The finite checks establish identities and authenticate literal equations,
degrees, counts, restored rows and fixed integers. They do not prove an
unbounded theorem from finite samples, search for a ZFC contradiction,
or rebuild the full transitive Lean/Rocq development. The separate
mathematical reviews in this package are separate AI-assisted reviews
within this work, not external human refereeing.

The inherited elementary Pell tools, rational integrality descent,
factorial-free selector, diagonal growth, strict-horizon change,
fixed-base power pooling and common-quadratic-direction projection
are credited to their retained sources. The new compact carriers,
smaller-modulus refinement, affine radix choice and exact compositions
are the contributions submitted for review. Improvement over this
pinned comparison set does not establish literature-wide priority.

% BEGIN reproduction.tex
\subsection{Reproducing and evaluating the deliverables}

The top-level \file{README.md} maps each selected equation to its exact
source and explains the component layouts. From the archive root,
run
\begin{verbatim}
python3 -B verify_release.py
python3 -B -O verify_release.py
\end{verbatim}
These read-only checks require Python 3.10 or newer and its standard
library. The optional \file{--receipt PATH} writes a fresh report
outside the sealed artifacts. The checker authenticates the complete
manifest and the literal selected equations before replaying the
component audits. It does not search for solutions.

The script \file{evaluate.py} evaluates a selected direct residual list
or compact arithmetic program at a supplied integer tuple. A modulus
can be supplied for manageable arithmetic with the compact constants.
Such modular evaluations are not tests of consistency. Exact zero
verification would require the complete candidate tuple in the declared
domain.

The human-readable principal residuals are generated from the canonical
equation JSON by \file{code/render_residuals.py}. To rebuild the article,
use a standard TeX installation:
\begin{verbatim}
cd report
pdflatex -interaction=nonstopmode -halt-on-error report.tex
pdflatex -interaction=nonstopmode -halt-on-error report.tex
\end{verbatim}
The supplied PDF was compiled and visually checked. Its source comprises
\file{report.tex} and the included \file{.tex} fragments. A flattened
\file{report_standalone.tex} is also supplied for convenient integration.
PDF metadata and byte layout may differ across TeX installations;
the literal polynomial data define the mathematical objects.

\subsection{Integration and the remaining formalization task}

The archive is an additive research submission. Existing repository
equations need not be overwritten. The source ledger pins the compared
commit, all twelve incoming archives, the immediate parents and the
retained licenses. Intermediate equations are kept to make each
restoration and source change reviewable; the selected outputs are
identified separately.

The next formalization task is to encode the new residue-carrier theorem,
its $q^4$ refinement, the rational-carrier application, and the concrete
direct composition against the existing arithmetic and logical interfaces.
The finite certificate checkers substantially narrow the implementation
surface, but they do not discharge that unbounded formalization task.

% END reproduction.tex


\appendix
\section{The full compact coefficient recipe}\label{app:compact}
% BEGIN compact_specification.tex
This is the complete finite indexed943/radix67 recipe retained from the
native compiler. In the notation below, $K_0=K$. The value $M_F$ is
already shifted by $B-1$ exactly once. No controller prime, Diophantine
witness, or consistency-dependent choice occurs in these definitions.
All displayed finite powers and sums are evaluated when specifying
the fixed coefficients; they do not add unknowns to an equation.

The canonical unsigned little-endian payloads have the following sizes:
\begin{center}\small
\begin{tabular}{@{}lrr@{}}
\toprule
Constant & Bits & Raw bytes\\
\midrule
$\beta$ &590,069&73,759\\
$K_0$ &592,884&74,111\\
$M_C$ &590,069&73,759\\
$M_F$ &590,070&73,759\\
\midrule
Total&&295,388\\
\bottomrule
\end{tabular}
\end{center}
The package supplies both these complete bytes and an independent
reconstruction script. The largest fixed input is 592,884 bits;
coefficients of an expanded high-degree output can be larger.
The exact equations are the straight-line programs with these fixed
integer leaves, not merely a family with four unspecified parameters.

\begingroup\small
% BEGIN constant_recipe.tex
\textbf{The four fixed integers.}
Every index, set, sum, product and floor below defines a constant, not
another unknown. Put
\[
\rho=2^{67},\quad B=\rho^{8807},\quad\beta=B-1,\quad
G_m(x)=\sum_{j=0}^{m-1}x^j.
\]
Define
\[
\begin{split}
\mathcal E={}&\rg0{942}\cup\rg{983}{992}\cup\rg{1003}{1011}\\
&\cup\rg{1013}{1031}\cup\{1063,1064,1095,2195,3296\}.
\end{split}
\]
Then
\begin{align*}
M_C={}&B-1-\sum_{e\in\mathcal E\setminus\{1\}}\rho^e-2\rho^{1063},\\
M_F={}&B-1+4+(2^{51}-2)(\rho^{4393}+\rho^{5495})\\
&+\rho^{1013}G_{10}(\rho)+\rho^{1065}G_9(\rho)
  +\rho^{2195}+\rho^{3296}\\
&+(2^{65}-1)\rho^{8792}G_{13}(\rho)+(2^{19}-1)\rho^{8805},\\
K={}&\rho^{30}+\rho^{62}+\rho^{1100}+\rho^{2201}+\rho^{8849}\\
&+(\rho^{4393}+\rho^{5495})\Phi(1/\rho),
\end{align*}
where
\[
\begin{split}
\Phi(x)={}&(1+2^{51})G_{943}(x)
 -2^{10}\sum_{i=0}^{942}v(t_i)x^i\\
&+\sum_{(b,e,m)\in\mathcal C}2^b x^e G_m(2x),\\
\mathcal C={}&\{(10,1023,8),(18,1013,9),(27,1003,8),\\
&(35,983,9),(44,1031,1),(45,1022,1),(46,1011,1),\\
&(47,992,1),(48,1095,1),(49,2195,1),(50,3296,1)\}.
\end{split}
\]
The rational intermediate in K cancels exactly; all four results are positive integers.

\textbf{The 943 indexed fixed tiles.}
Use the actions $u_j,v_j,d_j$ decoded below. Let
\[
\begin{gathered}
h(q,b)=2(q+1)+b,\quad
\mathcal R=\{u_j:d_j=1,\ 0\le j<430\},\\
\mathcal L=\{u_j:d_j=0,\ 0\le j<430\},\\
r_q=6+|\mathcal R\cap\rg0{q-1}|,\quad
l_q=205+|\mathcal L\cap\rg0{q-1}|.
\end{gathered}
\]
An interval ending at $-1$ is empty. The table defines $t_i$ for
$0\le i\le942$; all index ranges are fixed:
\[
\setlength{\arraycolsep}{1.4pt}\begin{array}{c|l}
\text{index }i&\text{tile }t_i=(w,n,e,s)\\\hline
0,1&(2,435,3,1),\ (5,0,0,0)\\
j+2\ (j=2q+b<430)&((1-d_j)l_{u_j},h(q,b),d_jr_{u_j},v_j)\\
432+b\ (b=0,1)&(0,h(215,b),0,h(215,b))\\
434+b\ (b=0,1)&(0,b,0,b)\\
424+2r_q+b\ (q\in\mathcal R,b=0,1)&(r_q,b,0,h(q,b))\\
424+2l_q+b\ (q\in\mathcal L,b=0,1)&(0,b,l_q,h(q,b))\\
924+j\ (0\le j<4)&(0,j+430\fl{j/2},0,438)\\
928&(2,434,2,0)\\
929&(3,436,4,4)\\
930&(4,436,4,0)\\
497+t\ (t=434,435,436)&(1,438,1,t)\\
\substack{934+3a+b-6\fl{a/4}\\a\in\{0,1,4\},\ b=0,1,2}&(a,437,b,437)
\end{array}
\]
For $t=(w,n,e,s)$ put
\[
v(t)=w+2^8n+2^{17}e+2^{25}s+127\,2^{34}.
\]

\textbf{A closed decoder for the 430 actions.}
Put $D_b(A,k)=\fl{A/b^k}\bmod b$, where
$a\bmod b=a-b\fl{a/b}$. For $0\le q<215$, $b=0,1$, $j=2q+b$:
\[
\begin{gathered}
N=\mathcal A\bmod2^{636},\quad
T_0=\fl{\mathcal A/2^{636}}\bmod2^{1675},\\
E_0=\fl{\mathcal A/2^{2311}},\\
\nu_j=D_2(N,j),\quad \sigma_j=\sum_{k=0}^{j-1}\nu_k,\\
m_q=D_2(N,430+q),\quad \tau_q=\sum_{k=0}^{q-1}m_k,\\
F_q=14(1-m_q)+m_qD_{16}(E_0,\tau_q),\\
u_j=\nu_j(1+\sigma_j)+(1-\nu_j)D_{216}(T_0,j-\sigma_j),\\
v_j=D_2(F_q,2b),\quad d_j=D_2(F_q,2b+1).
\end{gathered}
\]
Empty sums are zero. Here $(v_0,d_0,u_0)=(1,1,1)$; the constructor tile
uses this actual first target, including its state renaming.

\textbf{The integer $\mathcal A$.}
The alphabet, in order of digit values $0,\ldots,61$, is:\par
{\fontsize{9}{11}\selectfont\ttfamily 0123456789ABCDEFGHIJKLMNOPQRSTUVWXYZabcdefghijklmnopqrstuvwxyz}.\par
Concatenate the following 436 characters as digits $e_0,\ldots,e_{435}$;
put $\mathcal A=\sum_{j=0}^{435}e_j62^{435-j}$.\par
\begingroup\fontsize{9}{11}\selectfont\ttfamily
BxEnhvy6NipzkWmROo0GsA3uK3twWLxscLq20UTumy5Cdu461LwD6g7RDNzF8Y\par
iYDFTpB8ehAypd7d6OsOL4APamnJeNXQypUezCmW9n4oXP6XfgmiWuAWZ8ZJ7U\par
OP7gu4kVdlGZqZHULQMar1dXNeJW4i7oNVI8usp8BlovjjZbwPCovZI328fRUq\par
4q92R81s5GjWQaQfGhkq9UYMK3rRMaBmsOIBkJqc2cj2VRIZ6mjzkBIOkW9SjZ\par
FzeA6J5u4j3sBL3SdgwXzM9mHvQivu8EJ1TBTpRhHDqZIbmYgwQ8VqTAFotl6F\par
We1YC554GKrVHsA8oUGKSA6F6k5KREBSD0Tgd7bFALY84LLbpw574BQyvdoHrl\par
0Oe35uwDf15DpT6VJWZBjSbHPpPAINYTte0eUxnTFTKBlGj8Am6sb1gLldh8WZ\par
8P\par
\endgroup

% END constant_recipe.tex

\endgroup

The indexed tile recipe is calibrated by the retained native-language
proofs and their K215/refutation interface. The matching masks, shifted
port, action decoder and entire tile ordering are preserved together.
The separate newer radix47 tuple is not substituted into these equations.

\subsection*{An explicit unrestricted-integer companion}

For readers who require an unrestricted integer domain even on the
small-variable branch, replace each positive coordinate $u$ of
$P_{12,74}$ by
\[
 u=1+a^2+b^2+c(c+1),\qquad a,b,c\in\mathbb Z.
\]
Every resulting value is positive. Conversely, for any $u\ge1$,
the three-square theorem applied to $4(u-1)+1$ gives exactly one
odd square and two even squares. Dividing by four after subtracting
the odd square's constant term proves the displayed representation.
The substitution is therefore surjective onto the positive integers.

It gives a fully explicit equation in 36 unrestricted integer variables,
of exact degree 876, evaluated in 158 operations under the same model.
Each substitution costs seven operations; the degree doubles. The
distinguished degree-438 monomial survives by selecting only the $a^2$
term in every substituted coordinate, which proves exact attainment.
This is a standard domain conversion and is not claimed as a new
variable-count record. The literal lifted graph is included with the
selected equations and verified by the release checker.

% END compact_specification.tex


\section{Every quadratic residual of the principal equation}\label{app:rows}

All coordinates range independently over $\Z$. The equation is exactly
the sum of the squares of the following \DirectRows\ residuals.
Coordinate subscripts are indices, not sign restrictions.

% BEGIN residuals.tex
% Generated by code/render_residuals.py; edit the JSON source, not this file.
% Exact source SHA256: 6537b23ca3ab459d4b8215ea5c3c54d3d5880285eac86be4199acf2de2e7aed4
% Ten-point math; conservative 410-pt right-hand-side wrapping.
% Starred line breaks keep each individual residual together.
\begingroup
\fontsize{10}{12}\selectfont
\allowdisplaybreaks[4]
\setlength{\jot}{2pt}
\begin{align*}
% BEGIN_RESIDUAL 0
R_{0} ={}& x_{0} - x_{0}x_{2} - x_{2} + x_{7} \\[4pt]
% END_RESIDUAL 0
% BEGIN_RESIDUAL 1
R_{1} ={}& -1 + x_{1} - x_{2}x_{5} - x_{3} \\[4pt]
% END_RESIDUAL 1
% BEGIN_RESIDUAL 2
R_{2} ={}& -1 + x_{1} - x_{6}x_{7} \\[4pt]
% END_RESIDUAL 2
% BEGIN_RESIDUAL 3
R_{3} ={}& 1 + x_{3} - x_{4}x_{18} + x_{16} \\[4pt]
% END_RESIDUAL 3
% BEGIN_RESIDUAL 4
R_{4} ={}& -1 - x_{8}^{2} - 2x_{8}x_{9} - x_{9}^{2} + x_{19} \\[4pt]
% END_RESIDUAL 4
% BEGIN_RESIDUAL 5
R_{5} ={}& 7 - 8x_{8} + 2x_{8}^{2} - 4x_{8}x_{9} - 8x_{8}x_{19} - 16x_{9} - 10x_{9}^{2} - 16x_{9}x_{19} - 2x_{19} + 12x_{19}^{2} \\*
&{} + x_{20} - x_{36} \\[4pt]
% END_RESIDUAL 5
% BEGIN_RESIDUAL 6
R_{6} ={}& -2 + 2x_{8} - x_{8}^{2} - 4x_{8}x_{20} - 4x_{20} - 4x_{20}^{2} + x_{21} \\[4pt]
% END_RESIDUAL 6
% BEGIN_RESIDUAL 7
R_{7} ={}& -24x_{8}x_{21} + 24x_{8}x_{33} + 2x_{21} - 4x_{21}^{2} + x_{22} - 2x_{33} + 4x_{33}^{2} - x_{34} \\[4pt]
% END_RESIDUAL 7
% BEGIN_RESIDUAL 8
R_{8} ={}& -2 - x_{9} - x_{9}^{2} - 2x_{9}x_{10} - x_{10} - x_{10}^{2} + x_{23} \\[4pt]
% END_RESIDUAL 8
% BEGIN_RESIDUAL 9
R_{9} ={}& -1 + x_{8} - x_{8}^{2} - 2x_{8}x_{10} + x_{10} - x_{10}^{2} + x_{24} \\[4pt]
% END_RESIDUAL 9
% BEGIN_RESIDUAL 10
R_{10} ={}& 6 + x_{0} + x_{8} + 11x_{8}^{2} - 28x_{8}x_{9} - 4x_{8}x_{10} - 24x_{8}x_{19} - 4x_{8}x_{23} + 8x_{8}x_{24} \\*
&{} + 24x_{8}x_{37} + 5x_{9} + 11x_{9}^{2} - 14x_{9}x_{10} - 16x_{9}x_{19} + 8x_{9}x_{23} - 8x_{9}x_{24} - 3x_{10} \\*
&{} + x_{10}^{2} - 24x_{10}x_{19} - 8x_{10}x_{24} + 16x_{10}x_{37} - 16x_{19} - 4x_{19}^{2} - 8x_{19}x_{24} + 3x_{23} \\*
&{} - 4x_{23}^{2} - 8x_{23}x_{24} + 2x_{24} - 4x_{24}^{2} + 8x_{24}x_{37} + x_{25} - 4x_{25}^{2} + x_{26} + x_{27} \\*
&{} - 4x_{37} + 4x_{37}^{2} - x_{38} - 3x_{43} - 2x_{188}^{2} \\[4pt]
% END_RESIDUAL 10
% BEGIN_RESIDUAL 11
R_{11} ={}& -2 - x_{0} + 8x_{8} - 2x_{8}^{2} + 8x_{8}x_{10} - 8x_{8}x_{23} - 4x_{10} + 4x_{10}^{2} - 8x_{10}x_{23} + 4x_{23}^{2} \\*
&{} - 2x_{24} - 4x_{25}^{2} + x_{26} - 4x_{31} - 4x_{31}^{2} + x_{32} - 3x_{43} + 2x_{188}^{2} \\[4pt]
% END_RESIDUAL 11
% BEGIN_RESIDUAL 12
R_{12} ={}& -2 - 8x_{8} + 8x_{8}x_{23} + 8x_{8}x_{24} - 8x_{9} + 8x_{9}x_{23} + 8x_{9}x_{24} - 4x_{19} + 8x_{23} - 4x_{23}^{2} \\*
&{} - 8x_{23}x_{24} + 8x_{24} - 4x_{24}^{2} + x_{27} \\[4pt]
% END_RESIDUAL 12
% BEGIN_RESIDUAL 13
R_{13} ={}& -2 - 6x_{8} + 9x_{8}^{2} + 10x_{8}x_{9} - 8x_{8}x_{10} + 8x_{8}x_{19} - 16x_{8}x_{24} - 8x_{8}x_{27} + 4x_{9} \\*
&{} - 19x_{9}^{2} - 16x_{9}x_{19} - 32x_{9}x_{24} - 16x_{9}x_{27} - x_{10} + x_{10}^{2} + 8x_{10}x_{24} + 4x_{10}x_{27} \\*
&{} + 5x_{19} - 4x_{19}^{2} - 16x_{19}x_{24} - 8x_{19}x_{27} + 8x_{24} + 4x_{27} + x_{28} - x_{42} \\[4pt]
% END_RESIDUAL 13
% BEGIN_RESIDUAL 14
R_{14} ={}& 2 - x_{0} - 4x_{8} - 4x_{8}^{2} + 9x_{8}x_{9} - 9x_{8}x_{10} + 24x_{8}x_{19} - 4x_{8}x_{23} - 8x_{8}x_{24} + 4x_{8}x_{52} \\*
&{} - 10x_{9} - 3x_{9}^{2} + 13x_{9}x_{10} + 16x_{9}x_{19} + 8x_{9}x_{24} - 20x_{10} - 16x_{10}^{2} + 24x_{10}x_{19} \\*
&{} - 16x_{10}x_{23} - 24x_{10}x_{24} + 8x_{10}x_{52} + 16x_{19} + 4x_{19}^{2} + 8x_{19}x_{24} - 9x_{23} + 8x_{23}^{2} \\*
&{} - 14x_{24} - 4x_{24}^{2} - x_{25} - 4x_{25}^{2} + x_{26} + 2x_{29} + 4x_{29}^{2} - x_{30} + 16x_{31} + 16x_{31}^{2} \\*
&{} - 4x_{32} - 3x_{43} + 4x_{52} + 4x_{52}^{2} - x_{53} - 5x_{71} + 2x_{188}^{2} \\[4pt]
% END_RESIDUAL 14
% BEGIN_RESIDUAL 15
R_{15} ={}& -1 + 3x_{0} + 6x_{8} + x_{8}^{2} - 4x_{8}x_{9} + 2x_{8}x_{10} + 4x_{8}x_{23} + 16x_{8}x_{24} + 4x_{8}x_{39} - 3x_{9} \\*
&{} - x_{9}^{2} - 4x_{9}x_{10} - 16x_{9}x_{23} + 4x_{9}x_{24} + 7x_{10} - 6x_{10}^{2} - 8x_{10}x_{23} + 8x_{10}x_{24} \\*
&{} + 8x_{10}x_{39} - 12x_{10}x_{43} - 2x_{23} - 12x_{23}^{2} + 8x_{23}x_{24} + 16x_{23}x_{43} - x_{24} + 8x_{24}^{2} \\*
&{} + 4x_{25}^{2} - x_{26} - x_{29} - 4x_{29}^{2} + x_{30} - 4x_{31} - 4x_{31}^{2} + x_{32} - 4x_{39} + 4x_{39}^{2} \\*
&{} - x_{40} + 9x_{43} + 4x_{43}^{2} - x_{44} + 4x_{71} - 6x_{188}^{2} \\[4pt]
% END_RESIDUAL 15
% BEGIN_RESIDUAL 16
R_{16} ={}& 4 - 4x_{8} + 5x_{8}^{2} + 5x_{8}x_{9} - 3x_{8}x_{10} - 8x_{8}x_{24} - 8x_{8}x_{59} - 22x_{9} - 7x_{9}^{2} \\*
&{} + 13x_{9}x_{10} + 8x_{9}x_{19} + 16x_{9}x_{23} - 8x_{9}x_{24} - 16x_{9}x_{59} + 2x_{10} + 4x_{10}x_{57} + 8x_{19} \\*
&{} - 8x_{19}x_{23} + 8x_{19}x_{59} + x_{23} + 8x_{23}x_{24} - 6x_{24} + 4x_{24}^{2} - x_{27} + x_{31} + 4x_{57}^{2} \\*
&{} - x_{58} - 4x_{59} + 4x_{59}^{2} + 3x_{71} \\[4pt]
% END_RESIDUAL 16
% BEGIN_RESIDUAL 17
R_{17} ={}& -3 - 6x_{8} - 4x_{10} - 2x_{24} - 4x_{31} - 4x_{31}^{2} + x_{32} \\[4pt]
% END_RESIDUAL 17
% BEGIN_RESIDUAL 18
R_{18} ={}& 2 + x_{0} - 5x_{1} + 5x_{6}x_{7} + 11x_{8} + 14x_{8}^{2} + 18x_{8}x_{9} - 10x_{8}x_{10} + 16x_{8}x_{20} \\*
&{} + 8x_{8}x_{23} + 5x_{9} - 31x_{9}^{2} + 8x_{9}x_{10} - 24x_{9}x_{19} + 64x_{9}x_{20} + 24x_{9}x_{23} - 2x_{10} \\*
&{} + 6x_{10}^{2} + 4x_{10}x_{19} + 17x_{19} - 24x_{19}^{2} + 32x_{19}x_{20} - 8x_{19}x_{23} + 2x_{20} - 16x_{20}^{2} \\*
&{} + 3x_{21} - 2x_{22} + 6x_{23} - 4x_{23}^{2} - 4x_{24} - 2x_{26} - 2x_{28} - 2x_{30} - 11x_{31} - 12x_{31}^{2} \\*
&{} + x_{32} + x_{33} - 4x_{34} - 2x_{35} + x_{36} - 2x_{38} - 2x_{41} - x_{59} - 4x_{60} - 2x_{61} \\*
&{} - 4x_{185}x_{198} - 2x_{188}^{2} + x_{192} \\[4pt]
% END_RESIDUAL 18
% BEGIN_RESIDUAL 19
R_{19} ={}& 8 + 2x_{8} - 29x_{8}^{2} + 6x_{8}x_{9} + 8x_{8}x_{10} - 24x_{8}x_{33} + 3x_{9}^{2} - 4x_{10} + 4x_{10}^{2} - 3x_{19} \\*
&{} - 4x_{24} + 2x_{33} - 4x_{33}^{2} + x_{34} \\[4pt]
% END_RESIDUAL 19
% BEGIN_RESIDUAL 20
R_{20} ={}& -x_{9}^{2} - 4x_{9}x_{34} - 4x_{34}^{2} + x_{35} \\[4pt]
% END_RESIDUAL 20
% BEGIN_RESIDUAL 21
R_{21} ={}& -6 - 2x_{8} - 4x_{9} - 14x_{19} - 16x_{19}^{2} + x_{36} \\[4pt]
% END_RESIDUAL 21
% BEGIN_RESIDUAL 22
R_{22} ={}& 3 + x_{0} + 18x_{9} - 20x_{9}^{2} - 12x_{9}x_{10} - 24x_{9}x_{23} - 24x_{9}x_{36} + 4x_{9}x_{44} + 4x_{9}x_{47} \\*
&{} + 2x_{10} - 2x_{10}^{2} - 16x_{10}x_{36} + 8x_{10}x_{44} + 8x_{10}x_{47} + 4x_{23} - 12x_{23}^{2} - 8x_{23}x_{36} \\*
&{} + 4x_{36} - 4x_{36}^{2} + x_{37} + 2x_{43} + 4x_{44}^{2} - x_{45} + 4x_{47}^{2} - x_{48} - 2x_{188}^{2} \\[4pt]
% END_RESIDUAL 22
% BEGIN_RESIDUAL 23
R_{23} ={}& 4 - 20x_{8}^{2} - 4x_{8}x_{10} - 16x_{8}x_{24} - 24x_{8}x_{37} - 8x_{9} + 8x_{9}x_{23} + 8x_{9}x_{24} - 10x_{10} \\*
&{} + 6x_{10}^{2} - 8x_{10}x_{23} - 16x_{10}x_{24} - 16x_{10}x_{37} - 4x_{19} + 8x_{23} - 8x_{23}x_{24} - 6x_{24} \\*
&{} - 8x_{24}^{2} - 8x_{24}x_{37} - 4x_{25}^{2} + x_{26} + x_{27} + 4x_{31} + 4x_{31}^{2} - x_{32} + 4x_{37} - 4x_{37}^{2} \\*
&{} + x_{38} - 3x_{43} \\[4pt]
% END_RESIDUAL 23
% BEGIN_RESIDUAL 24
R_{24} ={}& -x_{10}^{2} - 2x_{10}x_{11} - x_{11}^{2} + x_{39} \\[4pt]
% END_RESIDUAL 24
% BEGIN_RESIDUAL 25
R_{25} ={}& -2 - x_{0} - 2x_{8} + 3x_{8}^{2} + 4x_{8}x_{9} - 4x_{8}x_{39} + x_{9} + 5x_{9}^{2} + 6x_{9}x_{10} - 15x_{10} \\*
&{} + 5x_{10}^{2} - 8x_{10}x_{23} + 16x_{10}x_{24} - 8x_{10}x_{39} - 5x_{23} + 8x_{23}^{2} - 8x_{23}x_{24} + 20x_{24} \\*
&{} - 4x_{24}^{2} + 4x_{39} - 4x_{39}^{2} + x_{40} - 3x_{43} - 4x_{71} + 2x_{188}^{2} \\[4pt]
% END_RESIDUAL 25
% BEGIN_RESIDUAL 26
R_{26} ={}& 3 + x_{8} + x_{8}^{2} - 2x_{8}x_{10} - 8x_{8}x_{39} + x_{10} - 5x_{10}^{2} - 16x_{10}x_{39} - 3x_{24} - 16x_{24}^{2} \\*
&{} - 32x_{24}x_{39} - 16x_{24}x_{40} + 4x_{39} - 24x_{39}^{2} - 16x_{39}x_{40} - 4x_{40}^{2} + x_{41} \\[4pt]
% END_RESIDUAL 26
% BEGIN_RESIDUAL 27
R_{27} ={}& -16x_{8}^{2} + 8x_{8}x_{10} + 32x_{8}x_{24} + 16x_{8}x_{27} + x_{10} - x_{10}^{2} - 8x_{10}x_{24} - 4x_{10}x_{27} \\*
&{} - 16x_{24}^{2} - 16x_{24}x_{27} - 4x_{27}^{2} + x_{42} \\[4pt]
% END_RESIDUAL 27
% BEGIN_RESIDUAL 28
R_{28} ={}& -1 - 8x_{10} - 4x_{10}^{2} + 8x_{10}x_{23} + 8x_{23} - 4x_{23}^{2} + x_{43} \\[4pt]
% END_RESIDUAL 28
% BEGIN_RESIDUAL 29
R_{29} ={}& -2 - x_{0} + 8x_{8} - 2x_{8}^{2} + 8x_{8}x_{10} - 8x_{8}x_{23} + 4x_{10} - x_{10}^{2} + 8x_{10}x_{23} + 12x_{10}x_{43} \\*
&{} - 8x_{23} - 8x_{23}^{2} - 16x_{23}x_{43} - 2x_{24} - 4x_{25}^{2} + x_{26} - 4x_{31} - 4x_{31}^{2} + x_{32} - 4x_{43} \\*
&{} - 4x_{43}^{2} + x_{44} + 2x_{188}^{2} \\[4pt]
% END_RESIDUAL 29
% BEGIN_RESIDUAL 30
R_{30} ={}& 1 + x_{9} - x_{9}^{2} - 4x_{9}x_{10} - 4x_{9}x_{44} - 4x_{10}^{2} - 8x_{10}x_{44} - 4x_{44}^{2} + x_{45} \\[4pt]
% END_RESIDUAL 30
% BEGIN_RESIDUAL 31
R_{31} ={}& -1 - 4x_{45} - 4x_{45}^{2} + x_{46} \\[4pt]
% END_RESIDUAL 31
% BEGIN_RESIDUAL 32
R_{32} ={}& 1 - x_{10}^{2} - 8x_{10}x_{45} - 4x_{10}x_{46} + 8x_{45} - 8x_{45}^{2} - 16x_{45}x_{46} - 2x_{46} - 4x_{46}^{2} + x_{47} \\[4pt]
% END_RESIDUAL 32
% BEGIN_RESIDUAL 33
R_{33} ={}& 2 + x_{9} - x_{9}^{2} - 4x_{9}x_{10} - 4x_{9}x_{47} - 4x_{10}^{2} - 8x_{10}x_{47} - 4x_{47}^{2} + x_{48} \\[4pt]
% END_RESIDUAL 33
% BEGIN_RESIDUAL 34
R_{34} ={}& -4x_{42}^{2} - 8x_{42}x_{48} - 4x_{48}^{2} + x_{49} \\[4pt]
% END_RESIDUAL 34
% BEGIN_RESIDUAL 35
R_{35} ={}& -9 + 7x_{8} - x_{8}^{2} - 8x_{8}x_{42} - 4x_{8}x_{49} + 24x_{42} - 12x_{42}^{2} + 8x_{42}x_{48} - 16x_{42}x_{49} \\*
&{} + 4x_{48}^{2} + 11x_{49} - 4x_{49}^{2} + x_{50} \\[4pt]
% END_RESIDUAL 35
% BEGIN_RESIDUAL 36
R_{36} ={}& -1 - 4x_{50} - 4x_{50}^{2} + x_{51} \\[4pt]
% END_RESIDUAL 36
% BEGIN_RESIDUAL 37
R_{37} ={}& 2 - x_{10}^{2} - 8x_{10}x_{50} - 4x_{10}x_{51} + 8x_{50} - 8x_{50}^{2} - 16x_{50}x_{51} - 2x_{51} - 4x_{51}^{2} + x_{52} \\[4pt]
% END_RESIDUAL 37
% BEGIN_RESIDUAL 38
R_{38} ={}& -x_{8} - x_{8}^{2} - 4x_{8}x_{10} - 4x_{8}x_{52} - 4x_{10} - 4x_{10}^{2} - 8x_{10}x_{52} - 4x_{52} - 4x_{52}^{2} + x_{53} \\[4pt]
% END_RESIDUAL 38
% BEGIN_RESIDUAL 39
R_{39} ={}& -10 - 12x_{53} - 4x_{53}^{2} + x_{54} \\[4pt]
% END_RESIDUAL 39
% BEGIN_RESIDUAL 40
R_{40} ={}& 2 + 4x_{10} + 16x_{10}x_{53} + 8x_{10}x_{54} - 4x_{23} - 16x_{23}x_{53} - 8x_{23}x_{54} - x_{43} + 8x_{53} \\*
&{} - 16x_{53}^{2} - 16x_{53}x_{54} + 4x_{54} - 4x_{54}^{2} + x_{55} \\[4pt]
% END_RESIDUAL 40
% BEGIN_RESIDUAL 41
R_{41} ={}& -9 + x_{9} - x_{9}^{2} + 8x_{9}x_{10} - 8x_{9}x_{23} - 4x_{9}x_{55} - 7x_{10}^{2} + 8x_{10}x_{23} - 12x_{10}x_{43} \\*
&{} + 16x_{10}x_{55} + 16x_{23}x_{43} - 16x_{23}x_{55} + 4x_{43}^{2} - x_{44} - 12x_{48} - 4x_{48}^{2} - 4x_{55}^{2} \\*
&{} + x_{56} \\[4pt]
% END_RESIDUAL 41
% BEGIN_RESIDUAL 42
R_{42} ={}& 5 + x_{0} + 24x_{48} - 8x_{48}^{2} - 16x_{48}x_{56} - 2x_{56} - 4x_{56}^{2} + x_{57} - 2x_{188}^{2} \\[4pt]
% END_RESIDUAL 42
% BEGIN_RESIDUAL 43
R_{43} ={}& 1 - 8x_{8} - 4x_{8}x_{9} + 8x_{8}x_{59} - 12x_{9} - 4x_{9}x_{10} + 16x_{9}x_{59} - 2x_{10} - 4x_{10}^{2} \\*
&{} + 4x_{10}x_{19} - 4x_{10}x_{57} + 8x_{19} - 8x_{19}x_{59} - 4x_{57}^{2} + x_{58} + 3x_{59} - 4x_{59}^{2} + 4x_{71} \\[4pt]
% END_RESIDUAL 43
% BEGIN_RESIDUAL 44
R_{44} ={}& 6 + 4x_{8} + 2x_{8}^{2} + 8x_{8}x_{19} + 4x_{9} - 10x_{9}^{2} + 4x_{9}x_{10} + 16x_{9}x_{19} + 3x_{10} + 3x_{10}^{2} \\*
&{} - 4x_{10}x_{19} - 10x_{19} - 4x_{19}^{2} + x_{59} - 4x_{71} \\[4pt]
% END_RESIDUAL 44
% BEGIN_RESIDUAL 45
R_{45} ={}& 3 + x_{0} - 8x_{8} + 3x_{8}^{2} - 9x_{8}x_{9} - 5x_{8}x_{10} + 8x_{8}x_{19} + 16x_{8}x_{58} - 2x_{9} - 2x_{9}^{2} \\*
&{} - 3x_{9}x_{10} + 16x_{9}x_{19} + 32x_{9}x_{58} + x_{10} + 5x_{10}^{2} + 12x_{10}x_{19} + 8x_{10}x_{57} + x_{19} \\*
&{} - 4x_{19}^{2} - 16x_{19}x_{58} - 7x_{23} + 8x_{57}^{2} - 2x_{58} - 4x_{58}^{2} - 3x_{59} + x_{60} + 7x_{71} \\*
&{} - 2x_{188}^{2} \\[4pt]
% END_RESIDUAL 45
% BEGIN_RESIDUAL 46
R_{46} ={}& -2 - 4x_{60} - 4x_{60}^{2} + x_{61} \\[4pt]
% END_RESIDUAL 46
% BEGIN_RESIDUAL 47
R_{47} ={}& -x_{3}x_{13} + x_{62} \\[4pt]
% END_RESIDUAL 47
% BEGIN_RESIDUAL 48
R_{48} ={}& x_{12} - x_{13} - 2x_{22}x_{62} - x_{63} \\[4pt]
% END_RESIDUAL 48
% BEGIN_RESIDUAL 49
R_{49} ={}& x_{12} - x_{13} - 2x_{26}x_{62} - x_{64} \\[4pt]
% END_RESIDUAL 49
% BEGIN_RESIDUAL 50
R_{50} ={}& x_{12} - x_{13} - 2x_{28}x_{62} - 6x_{62} - x_{65} \\[4pt]
% END_RESIDUAL 50
% BEGIN_RESIDUAL 51
R_{51} ={}& x_{12} - x_{13} - 2x_{30}x_{62} - x_{66} \\[4pt]
% END_RESIDUAL 51
% BEGIN_RESIDUAL 52
R_{52} ={}& x_{12} - x_{13} - 2x_{32}x_{62} + 6x_{62} - x_{67} \\[4pt]
% END_RESIDUAL 52
% BEGIN_RESIDUAL 53
R_{53} ={}& x_{12} - x_{13} - 4x_{34}x_{62} - 2x_{35}x_{62} - x_{68} \\[4pt]
% END_RESIDUAL 53
% BEGIN_RESIDUAL 54
R_{54} ={}& x_{12} - x_{13} - 2x_{38}x_{62} - x_{69} \\[4pt]
% END_RESIDUAL 54
% BEGIN_RESIDUAL 55
R_{55} ={}& x_{12} - x_{13} - 2x_{41}x_{62} - 4x_{62} - x_{70} \\[4pt]
% END_RESIDUAL 55
% BEGIN_RESIDUAL 56
R_{56} ={}& -x_{8}x_{9} + x_{8}x_{10} + x_{9} + x_{9}x_{10} - x_{10} - x_{10}^{2} + x_{71} \\[4pt]
% END_RESIDUAL 56
% BEGIN_RESIDUAL 57
R_{57} ={}& 1 + 2x_{11} + x_{11}^{2} - x_{71}^{2} + x_{72} \\[4pt]
% END_RESIDUAL 57
% BEGIN_RESIDUAL 58
R_{58} ={}& -4x_{60}x_{62} - 2x_{61}x_{62} + 2x_{62} + x_{73} \\[4pt]
% END_RESIDUAL 58
% BEGIN_RESIDUAL 59
R_{59} ={}& -x_{12}^{2} + 2x_{12}x_{13} + 2x_{12}x_{73} - x_{13}^{2} - 2x_{13}x_{73} - x_{72}^{2} - x_{73}^{2} + x_{74} \\[4pt]
% END_RESIDUAL 59
% BEGIN_RESIDUAL 60
R_{60} ={}& -1 - x_{3}x_{9} + x_{75} \\[4pt]
% END_RESIDUAL 60
% BEGIN_RESIDUAL 61
R_{61} ={}& -x_{13}x_{75} + x_{77}x_{78} + x_{79} \\[4pt]
% END_RESIDUAL 61
% BEGIN_RESIDUAL 62
R_{62} ={}& -1 + x_{3} - x_{3}x_{8} + x_{76} \\[4pt]
% END_RESIDUAL 62
% BEGIN_RESIDUAL 63
R_{63} ={}& -1 + 2x_{3} + 2x_{3}x_{8} + 4x_{3}x_{9} - 2x_{3}x_{19} + x_{77} \\[4pt]
% END_RESIDUAL 63
% BEGIN_RESIDUAL 64
R_{64} ={}& -x_{12}x_{76} + x_{78} \\[4pt]
% END_RESIDUAL 64
% BEGIN_RESIDUAL 65
R_{65} ={}& -6x_{3}x_{13} + x_{12}x_{75} + 12x_{13} - 8x_{13}x_{75} - 4x_{13}x_{76} - x_{13}x_{77} + x_{80} \\[4pt]
% END_RESIDUAL 65
% BEGIN_RESIDUAL 66
R_{66} ={}& -x_{12}x_{63} + x_{13}x_{63} + x_{81} \\[4pt]
% END_RESIDUAL 66
% BEGIN_RESIDUAL 67
R_{67} ={}& -x_{64}x_{81} + x_{82} \\[4pt]
% END_RESIDUAL 67
% BEGIN_RESIDUAL 68
R_{68} ={}& -x_{65}x_{82} + x_{83} \\[4pt]
% END_RESIDUAL 68
% BEGIN_RESIDUAL 69
R_{69} ={}& -x_{66}x_{83} + x_{84} \\[4pt]
% END_RESIDUAL 69
% BEGIN_RESIDUAL 70
R_{70} ={}& -x_{67}x_{84} + x_{85} \\[4pt]
% END_RESIDUAL 70
% BEGIN_RESIDUAL 71
R_{71} ={}& -x_{68}x_{85} + x_{86} \\[4pt]
% END_RESIDUAL 71
% BEGIN_RESIDUAL 72
R_{72} ={}& -x_{69}x_{86} + x_{87} \\[4pt]
% END_RESIDUAL 72
% BEGIN_RESIDUAL 73
R_{73} ={}& -x_{70}x_{87} + x_{88} \\[4pt]
% END_RESIDUAL 73
% BEGIN_RESIDUAL 74
R_{74} ={}& -x_{79}x_{88} + x_{89} \\[4pt]
% END_RESIDUAL 74
% BEGIN_RESIDUAL 75
R_{75} ={}& -x_{80}x_{89} + x_{90} \\[4pt]
% END_RESIDUAL 75
% BEGIN_RESIDUAL 76
R_{76} ={}& -x_{90}^{2} + x_{91} \\[4pt]
% END_RESIDUAL 76
% BEGIN_RESIDUAL 77
R_{77} ={}& -x_{2} - x_{2}x_{16} + x_{16} + x_{92} \\[4pt]
% END_RESIDUAL 77
% BEGIN_RESIDUAL 78
R_{78} ={}& -x_{14}x_{92} + x_{93} \\[4pt]
% END_RESIDUAL 78
% BEGIN_RESIDUAL 79
R_{79} ={}& -2x_{1}^{2} + 2x_{1}x_{12} + 2x_{1}x_{13} + 2x_{1}x_{93} + 2x_{1}x_{94} - x_{12}^{2} - 2x_{12}x_{93} - x_{13}^{2} \\*
&{} - 2x_{13}x_{94} + x_{17}x_{95} - x_{74}x_{91} - x_{93}^{2} - x_{94}^{2} \\[4pt]
% END_RESIDUAL 79
% BEGIN_RESIDUAL 80
R_{80} ={}& -x_{2}x_{15} + x_{15} - x_{15}x_{92} + x_{94} \\[4pt]
% END_RESIDUAL 80
% BEGIN_RESIDUAL 81
R_{81} ={}& -3x_{98} - 6x_{98}x_{211} + x_{100} \\[4pt]
% END_RESIDUAL 81
% BEGIN_RESIDUAL 82
R_{82} ={}& -x_{4}x_{121} + x_{96}x_{203} - x_{99} - x_{100}x_{121} \\[4pt]
% END_RESIDUAL 82
% BEGIN_RESIDUAL 83
R_{83} ={}& x_{101}x_{203} - x_{102}x_{124} - x_{103} \\[4pt]
% END_RESIDUAL 83
% BEGIN_RESIDUAL 84
R_{84} ={}& x_{104}x_{203} - x_{105}x_{124} - x_{106} \\[4pt]
% END_RESIDUAL 84
% BEGIN_RESIDUAL 85
R_{85} ={}& x_{107}x_{203} - x_{108}x_{124} - x_{109} \\[4pt]
% END_RESIDUAL 85
% BEGIN_RESIDUAL 86
R_{86} ={}& x_{110}x_{203} - x_{111}x_{124} - x_{112} \\[4pt]
% END_RESIDUAL 86
% BEGIN_RESIDUAL 87
R_{87} ={}& x_{113}x_{204} - x_{115}x_{124} - x_{116} \\[4pt]
% END_RESIDUAL 87
% BEGIN_RESIDUAL 88
R_{88} ={}& x_{117}x_{204} - x_{118}x_{124} - x_{119} \\[4pt]
% END_RESIDUAL 88
% BEGIN_RESIDUAL 89
R_{89} ={}& -x_{17}x_{121} + x_{124} \\[4pt]
% END_RESIDUAL 89
% BEGIN_RESIDUAL 90
R_{90} ={}& x_{120} - x_{122}x_{124} - x_{123} \\[4pt]
% END_RESIDUAL 90
% BEGIN_RESIDUAL 91
R_{91} ={}& -x_{124}x_{126} + x_{125} - x_{127} \\[4pt]
% END_RESIDUAL 91
% BEGIN_RESIDUAL 92
R_{92} ={}& -x_{124}x_{129} + x_{128}x_{203} - x_{130} \\[4pt]
% END_RESIDUAL 92
% BEGIN_RESIDUAL 93
R_{93} ={}& -x_{124}x_{132} + x_{131}x_{203} - x_{133} \\[4pt]
% END_RESIDUAL 93
% BEGIN_RESIDUAL 94
R_{94} ={}& x_{134} - 2x_{134}x_{136} + x_{137} \\[4pt]
% END_RESIDUAL 94
% BEGIN_RESIDUAL 95
R_{95} ={}& -x_{134}x_{137} - x_{137} + x_{138} \\[4pt]
% END_RESIDUAL 95
% BEGIN_RESIDUAL 96
R_{96} ={}& 2x_{135} + x_{135}^{2} + x_{137}^{2} - x_{138}^{2} \\[4pt]
% END_RESIDUAL 96
% BEGIN_RESIDUAL 97
R_{97} ={}& -x_{135}x_{140} - x_{140} + x_{141} \\[4pt]
% END_RESIDUAL 97
% BEGIN_RESIDUAL 98
R_{98} ={}& 2x_{139} + x_{139}^{2} + x_{140}^{2} - x_{141}^{2} \\[4pt]
% END_RESIDUAL 98
% BEGIN_RESIDUAL 99
R_{99} ={}& 1 - 2x_{0} + x_{0}^{2} - 2x_{0}x_{135} + x_{142} \\[4pt]
% END_RESIDUAL 99
% BEGIN_RESIDUAL 100
R_{100} ={}& -1 - x_{0}x_{140} + x_{3} - x_{139} + x_{141} + x_{142}x_{143} \\[4pt]
% END_RESIDUAL 100
% BEGIN_RESIDUAL 101
R_{101} ={}& 1 - 2x_{4} + x_{4}^{2} - 2x_{4}x_{135} + x_{144} \\[4pt]
% END_RESIDUAL 101
% BEGIN_RESIDUAL 102
R_{102} ={}& -x_{4}x_{176} + x_{17} + x_{144}x_{145} - x_{175} + x_{177} \\[4pt]
% END_RESIDUAL 102
% BEGIN_RESIDUAL 103
R_{103} ={}& -1 + 2x_{135}x_{187} - x_{146} + 2x_{187} - x_{187}^{2} \\[4pt]
% END_RESIDUAL 103
% BEGIN_RESIDUAL 104
R_{104} ={}& 3 + x_{146}x_{147} - x_{175} - x_{176}x_{187} + x_{177} + 6x_{211} \\[4pt]
% END_RESIDUAL 104
% BEGIN_RESIDUAL 105
R_{105} ={}& -x_{135}x_{149} - x_{149} + x_{150} \\[4pt]
% END_RESIDUAL 105
% BEGIN_RESIDUAL 106
R_{106} ={}& 2x_{148} + x_{148}^{2} + x_{149}^{2} - x_{150}^{2} \\[4pt]
% END_RESIDUAL 106
% BEGIN_RESIDUAL 107
R_{107} ={}& 1 + x_{135} - x_{135}x_{152} - x_{152} + x_{153} \\[4pt]
% END_RESIDUAL 107
% BEGIN_RESIDUAL 108
R_{108} ={}& 1 + 2x_{151} + x_{151}^{2} - 2x_{152} + x_{152}^{2} - x_{153}^{2} \\[4pt]
% END_RESIDUAL 108
% BEGIN_RESIDUAL 109
R_{109} ={}& -x_{135}x_{155} - x_{155} + x_{156} \\[4pt]
% END_RESIDUAL 109
% BEGIN_RESIDUAL 110
R_{110} ={}& 2x_{154} + x_{154}^{2} + x_{155}^{2} - x_{156}^{2} \\[4pt]
% END_RESIDUAL 110
% BEGIN_RESIDUAL 111
R_{111} ={}& -x_{135}x_{158} - x_{158} + x_{159} \\[4pt]
% END_RESIDUAL 111
% BEGIN_RESIDUAL 112
R_{112} ={}& 2x_{157} + x_{157}^{2} + x_{158}^{2} - x_{159}^{2} \\[4pt]
% END_RESIDUAL 112
% BEGIN_RESIDUAL 113
R_{113} ={}& -x_{135}x_{161} - x_{161} + x_{162} \\[4pt]
% END_RESIDUAL 113
% BEGIN_RESIDUAL 114
R_{114} ={}& 2x_{160} + x_{160}^{2} + x_{161}^{2} - x_{162}^{2} \\[4pt]
% END_RESIDUAL 114
% BEGIN_RESIDUAL 115
R_{115} ={}& -x_{135}x_{164} - x_{164} + x_{165} \\[4pt]
% END_RESIDUAL 115
% BEGIN_RESIDUAL 116
R_{116} ={}& 2x_{163} + x_{163}^{2} + x_{164}^{2} - x_{165}^{2} \\[4pt]
% END_RESIDUAL 116
% BEGIN_RESIDUAL 117
R_{117} ={}& -x_{135}x_{167} - x_{167} + x_{168} \\[4pt]
% END_RESIDUAL 117
% BEGIN_RESIDUAL 118
R_{118} ={}& 2x_{166} + x_{166}^{2} + x_{167}^{2} - x_{168}^{2} \\[4pt]
% END_RESIDUAL 118
% BEGIN_RESIDUAL 119
R_{119} ={}& -x_{135}x_{170} - x_{170} + x_{171} \\[4pt]
% END_RESIDUAL 119
% BEGIN_RESIDUAL 120
R_{120} ={}& 2x_{169} + x_{169}^{2} + x_{170}^{2} - x_{171}^{2} \\[4pt]
% END_RESIDUAL 120
% BEGIN_RESIDUAL 121
R_{121} ={}& -x_{135}x_{173} - x_{173} + x_{174} \\[4pt]
% END_RESIDUAL 121
% BEGIN_RESIDUAL 122
R_{122} ={}& 2x_{172} + x_{172}^{2} + x_{173}^{2} - x_{174}^{2} \\[4pt]
% END_RESIDUAL 122
% BEGIN_RESIDUAL 123
R_{123} ={}& 1 - x_{135}x_{176} - x_{176} + x_{177} \\[4pt]
% END_RESIDUAL 123
% BEGIN_RESIDUAL 124
R_{124} ={}& -1 + 2x_{175} + x_{175}^{2} + x_{176}^{2} - 2x_{177} - x_{177}^{2} \\[4pt]
% END_RESIDUAL 124
% BEGIN_RESIDUAL 125
R_{125} ={}& -x_{135}x_{179} - x_{179} + x_{180} \\[4pt]
% END_RESIDUAL 125
% BEGIN_RESIDUAL 126
R_{126} ={}& 2x_{178} + x_{178}^{2} + x_{179}^{2} - x_{180}^{2} \\[4pt]
% END_RESIDUAL 126
% BEGIN_RESIDUAL 127
R_{127} ={}& -x_{135}x_{182} - x_{182} + x_{183} \\[4pt]
% END_RESIDUAL 127
% BEGIN_RESIDUAL 128
R_{128} ={}& 2x_{181} + x_{181}^{2} + x_{182}^{2} - x_{183}^{2} \\[4pt]
% END_RESIDUAL 128
% BEGIN_RESIDUAL 129
R_{129} ={}& -x_{135}x_{185} - x_{185} + x_{186} \\[4pt]
% END_RESIDUAL 129
% BEGIN_RESIDUAL 130
R_{130} ={}& 2x_{184} + x_{184}^{2} + x_{185}^{2} - x_{186}^{2} \\[4pt]
% END_RESIDUAL 130
% BEGIN_RESIDUAL 131
R_{131} ={}& 4 - 2x_{0} - 4x_{0}^{2} - 2x_{0}x_{3} - x_{1}^{2} - x_{2}^{2} - 2x_{3}^{2} - x_{4}^{2} + 2x_{8} - x_{8}^{2} - x_{9}^{2} - x_{10}^{2} \\*
&{} - x_{11}^{2} - x_{12}^{2} - x_{13}^{2} - x_{14}^{2} - x_{15}^{2} - x_{16}^{2} - x_{17}^{2} - x_{96}^{2} - x_{97}^{2} - x_{101}^{2} \\*
&{} - x_{104}^{2} - x_{107}^{2} - x_{110}^{2} - x_{113}^{2} - x_{114}^{2} - x_{117}^{2} - x_{120}^{2} - x_{121}^{2} - x_{125}^{2} - x_{128}^{2} \\*
&{} - x_{131}^{2} + x_{134} + x_{140}^{2} + 12x_{146}x_{147} + x_{149}^{2} - 2x_{152} + x_{152}^{2} + x_{155}^{2} + x_{158}^{2} \\*
&{} + x_{161}^{2} + x_{164}^{2} + x_{167}^{2} + x_{170}^{2} + x_{173}^{2} - 2x_{175} + 5x_{175}^{2} + 6x_{176}^{2} - 12x_{176}x_{187} \\*
&{} + 2x_{177} - 5x_{177}^{2} + x_{179}^{2} + x_{182}^{2} - x_{185}^{2} - x_{187}^{2} + x_{205}^{2} + x_{206}^{2} + x_{207} \\*
&{} + x_{207}^{2} - 12x_{211} - 72x_{211}^{2} \\[4pt]
% END_RESIDUAL 131
% BEGIN_RESIDUAL 132
R_{132} ={}& 3x_{4} - 6x_{4}x_{8} - 6x_{4}x_{9} - 6x_{4}x_{10} - 6x_{4}x_{11} - 6x_{4}x_{12} - 6x_{4}x_{13} - 6x_{4}x_{14} - 6x_{4}x_{15} \\*
&{} - 6x_{4}x_{16} + x_{187} \\[4pt]
% END_RESIDUAL 132
% BEGIN_RESIDUAL 133
R_{133} ={}& -11 + x_{0} - 5x_{1} + 5x_{6}x_{7} - 2x_{188}^{2} \\[4pt]
% END_RESIDUAL 133
% BEGIN_RESIDUAL 134
R_{134} ={}& -x_{135}x_{190} - x_{190} + x_{199} \\[4pt]
% END_RESIDUAL 134
% BEGIN_RESIDUAL 135
R_{135} ={}& -1 + x_{189}^{2} + x_{190}^{2} - x_{199}^{2} \\[4pt]
% END_RESIDUAL 135
% BEGIN_RESIDUAL 136
R_{136} ={}& -x_{192} - x_{192}x_{193} + x_{200} \\[4pt]
% END_RESIDUAL 136
% BEGIN_RESIDUAL 137
R_{137} ={}& -1 + x_{191}^{2} + x_{192}^{2} - x_{200}^{2} \\[4pt]
% END_RESIDUAL 137
% BEGIN_RESIDUAL 138
R_{138} ={}& -x_{185}^{2} + x_{201} \\[4pt]
% END_RESIDUAL 138
% BEGIN_RESIDUAL 139
R_{139} ={}& x_{190} - 2x_{195}x_{201} - x_{201} \\[4pt]
% END_RESIDUAL 139
% BEGIN_RESIDUAL 140
R_{140} ={}& -4x_{190} - 8x_{190}x_{194} + x_{193} \\[4pt]
% END_RESIDUAL 140
% BEGIN_RESIDUAL 141
R_{141} ={}& -x_{135} - x_{189}x_{196} + x_{193} \\[4pt]
% END_RESIDUAL 141
% BEGIN_RESIDUAL 142
R_{142} ={}& -1 - x_{184} - x_{189}x_{197} + x_{191} \\[4pt]
% END_RESIDUAL 142
% BEGIN_RESIDUAL 143
R_{143} ={}& -x_{8} - x_{9} - 4x_{19} - 2x_{22} - 2x_{26} - 2x_{28} - 2x_{30} - 2x_{32} - 4x_{34} - 2x_{35} - 2x_{38} \\*
&{} - 2x_{41} - 4x_{60} - 2x_{61} - 4x_{185}x_{198} + x_{192} \\[4pt]
% END_RESIDUAL 143
% BEGIN_RESIDUAL 144
R_{144} ={}& -2 + 2x_{0}^{2} + 2x_{0}x_{3} - 2x_{0}x_{140} - 2x_{0}x_{149} + x_{1}^{2} - 2x_{1}x_{176} + x_{2}^{2} - 2x_{2}x_{167} \\*
&{} + x_{3}^{2} - 2x_{3}x_{149} + x_{8}^{2} - 2x_{8}x_{152} - 2x_{8}x_{185} + x_{9}^{2} - 2x_{9}x_{155} - 2x_{9}x_{185} + x_{10}^{2} \\*
&{} - 2x_{10}x_{158} + x_{11}^{2} - 2x_{11}x_{161} + x_{12}^{2} - 2x_{12}x_{164} + x_{13}^{2} - 2x_{13}x_{170} + x_{14}^{2} \\*
&{} - 2x_{14}x_{173} + x_{15}^{2} - 2x_{15}x_{179} + x_{16}^{2} - 2x_{16}x_{182} - 8x_{19}x_{185} - 4x_{22}x_{185} \\*
&{} - 4x_{26}x_{185} - 4x_{28}x_{185} - 4x_{30}x_{185} - 4x_{32}x_{185} - 8x_{34}x_{185} - 4x_{35}x_{185} - 4x_{38}x_{185} \\*
&{} - 4x_{41}x_{185} - 8x_{60}x_{185} - 4x_{61}x_{185} - x_{135}x_{202} + x_{139}^{2} + x_{148}^{2} + x_{151}^{2} + 2x_{152} \\*
&{} + x_{154}^{2} + x_{157}^{2} + x_{160}^{2} + x_{163}^{2} + x_{166}^{2} + x_{169}^{2} + x_{172}^{2} + x_{175}^{2} + x_{178}^{2} + x_{181}^{2} \\*
&{} + x_{184}^{2} + 2x_{185}^{2} - x_{205}^{2} - x_{206}^{2} - x_{207} - x_{207}^{2} \\[4pt]
% END_RESIDUAL 144
% BEGIN_RESIDUAL 145
R_{145} ={}& -x_{97}x_{203} + x_{121} \\[4pt]
% END_RESIDUAL 145
% BEGIN_RESIDUAL 146
R_{146} ={}& -x_{114}x_{204} + x_{121} \\[4pt]
% END_RESIDUAL 146
% BEGIN_RESIDUAL 147
R_{147} ={}& -1 + x_{2} - x_{3} - 2x_{3}x_{212} \\[4pt]
% END_RESIDUAL 147
% BEGIN_RESIDUAL 148
R_{148} ={}& 2 - x_{8}^{2} - 2x_{8}x_{9} - 8x_{8}x_{19} - 4x_{8}x_{22} - 4x_{8}x_{26} - 4x_{8}x_{28} - 4x_{8}x_{30} - 4x_{8}x_{32} \\*
&{} - 8x_{8}x_{34} - 4x_{8}x_{35} - 4x_{8}x_{38} - 4x_{8}x_{41} - 8x_{8}x_{60} - 4x_{8}x_{61} + 2x_{9}^{2} - 8x_{9}x_{19} \\*
&{} - 4x_{9}x_{22} - 4x_{9}x_{26} - 4x_{9}x_{28} - 4x_{9}x_{30} - 4x_{9}x_{32} + 4x_{9}x_{34} - 4x_{9}x_{35} - 4x_{9}x_{38} \\*
&{} - 4x_{9}x_{41} - 8x_{9}x_{60} - 4x_{9}x_{61} - 16x_{19}^{2} - 16x_{19}x_{22} - 16x_{19}x_{26} - 16x_{19}x_{28} \\*
&{} - 16x_{19}x_{30} - 16x_{19}x_{32} - 32x_{19}x_{34} - 16x_{19}x_{35} - 16x_{19}x_{38} - 16x_{19}x_{41} \\*
&{} - 32x_{19}x_{60} - 16x_{19}x_{61} - 4x_{22}^{2} - 8x_{22}x_{26} - 8x_{22}x_{28} - 8x_{22}x_{30} - 8x_{22}x_{32} \\*
&{} - 16x_{22}x_{34} - 8x_{22}x_{35} - 8x_{22}x_{38} - 8x_{22}x_{41} - 16x_{22}x_{60} - 8x_{22}x_{61} - 4x_{26}^{2} \\*
&{} - 8x_{26}x_{28} - 8x_{26}x_{30} - 8x_{26}x_{32} - 16x_{26}x_{34} - 8x_{26}x_{35} - 8x_{26}x_{38} - 8x_{26}x_{41} \\*
&{} - 16x_{26}x_{60} - 8x_{26}x_{61} - 4x_{28}^{2} - 8x_{28}x_{30} - 8x_{28}x_{32} - 16x_{28}x_{34} - 8x_{28}x_{35} \\*
&{} - 8x_{28}x_{38} - 8x_{28}x_{41} - 16x_{28}x_{60} - 8x_{28}x_{61} - 4x_{30}^{2} - 8x_{30}x_{32} - 16x_{30}x_{34} \\*
&{} - 8x_{30}x_{35} - 8x_{30}x_{38} - 8x_{30}x_{41} - 16x_{30}x_{60} - 8x_{30}x_{61} - 4x_{32}^{2} - 16x_{32}x_{34} \\*
&{} - 8x_{32}x_{35} - 8x_{32}x_{38} - 8x_{32}x_{41} - 16x_{32}x_{60} - 8x_{32}x_{61} - 4x_{34}^{2} - 16x_{34}x_{35} \\*
&{} - 16x_{34}x_{38} - 16x_{34}x_{41} - 32x_{34}x_{60} - 16x_{34}x_{61} - 3x_{35} - 4x_{35}^{2} - 8x_{35}x_{38} \\*
&{} - 8x_{35}x_{41} - 16x_{35}x_{60} - 8x_{35}x_{61} - 4x_{38}^{2} - 8x_{38}x_{41} - 16x_{38}x_{60} - 8x_{38}x_{61} \\*
&{} - 4x_{41}^{2} - 16x_{41}x_{60} - 8x_{41}x_{61} + 8x_{60} - 8x_{60}^{2} - 16x_{60}x_{61} - 2x_{61} - 4x_{61}^{2} \\*
&{} - x_{140}^{2} - x_{149}^{2} + 2x_{152} - x_{152}^{2} - x_{155}^{2} - x_{158}^{2} - x_{161}^{2} - x_{164}^{2} - x_{167}^{2} \\*
&{} - x_{170}^{2} - x_{173}^{2} - x_{176}^{2} - x_{179}^{2} - x_{182}^{2} + x_{185}^{2} - x_{205}^{2} - x_{206}^{2} - x_{207} - x_{207}^{2} \\[4pt]
% END_RESIDUAL 148
% BEGIN_RESIDUAL 149
R_{149} ={}& -1 - 4x_{99}^{2} - 4x_{103}^{2} - 4x_{106}^{2} - 4x_{109}^{2} - 4x_{112}^{2} - 4x_{116}^{2} - 4x_{119}^{2} + x_{121}^{2} \\*
&{} - 4x_{123}^{2} - 4x_{127}^{2} - 4x_{130}^{2} - 4x_{133}^{2} - x_{208}^{2} - x_{209}^{2} - x_{210} - x_{210}^{2} \\[4pt]
% END_RESIDUAL 149
% BEGIN_RESIDUAL 150
R_{150} ={}& -2 - 2x_{96}x_{149} + 2x_{101} - 2x_{101}x_{152} - 2x_{104}x_{155} - 2x_{107}x_{158} - 2x_{110}x_{161} \\*
&{} - 2x_{113}x_{164} - 2x_{117}x_{170} - 2x_{120}x_{173} - 2x_{125}x_{179} - 2x_{128}x_{182} - 2x_{131}x_{185} \\*
&{} - 2x_{135}x_{214} + 2x_{148}x_{149} + 2x_{149} - 2x_{151} + 2x_{151}x_{152} + 2x_{152} + 2x_{154}x_{155} \\*
&{} + 2x_{155} + 2x_{157}x_{158} + 2x_{158} + 2x_{160}x_{161} + 2x_{161} + 2x_{163}x_{164} + 2x_{164} \\*
&{} + 2x_{169}x_{170} + 2x_{170} + 2x_{172}x_{173} + 2x_{173} + 2x_{178}x_{179} + 2x_{179} + 2x_{181}x_{182} \\*
&{} + 2x_{182} + 2x_{184}x_{185} + 2x_{185} + 6x_{211}x_{214} + x_{213} + 2x_{214} \\[4pt]
% END_RESIDUAL 150
% BEGIN_RESIDUAL 151
R_{151} ={}& 10 - 2x_{96} + x_{96}^{2} - 2x_{96}x_{148} + 2x_{96}x_{150} - 2x_{101} + x_{101}^{2} - 2x_{101}x_{151} \\*
&{} + 2x_{101}x_{153} - 2x_{104} + x_{104}^{2} - 2x_{104}x_{154} + 2x_{104}x_{156} - 2x_{107} + x_{107}^{2} \\*
&{} - 2x_{107}x_{157} + 2x_{107}x_{159} - 2x_{110} + x_{110}^{2} - 2x_{110}x_{160} + 2x_{110}x_{162} - 2x_{113} \\*
&{} + x_{113}^{2} - 2x_{113}x_{163} + 2x_{113}x_{165} - 2x_{117} + x_{117}^{2} - 2x_{117}x_{169} + 2x_{117}x_{171} \\*
&{} - 2x_{120} + x_{120}^{2} - 2x_{120}x_{172} + 2x_{120}x_{174} - 2x_{125} + x_{125}^{2} - 2x_{125}x_{178} \\*
&{} + 2x_{125}x_{180} - 2x_{128} + x_{128}^{2} - 2x_{128}x_{181} + 2x_{128}x_{183} - 2x_{131} + x_{131}^{2} \\*
&{} - 2x_{131}x_{184} + 2x_{131}x_{186} + 2x_{148} + x_{148}^{2} - 2x_{148}x_{150} - x_{149}^{2} - 2x_{150} + x_{150}^{2} \\*
&{} + 2x_{151} + x_{151}^{2} - 2x_{151}x_{153} + 2x_{152} - x_{152}^{2} - 2x_{153} + x_{153}^{2} + 2x_{154} + x_{154}^{2} \\*
&{} - 2x_{154}x_{156} - x_{155}^{2} - 2x_{156} + x_{156}^{2} + 2x_{157} + x_{157}^{2} - 2x_{157}x_{159} - x_{158}^{2} \\*
&{} - 2x_{159} + x_{159}^{2} + 2x_{160} + x_{160}^{2} - 2x_{160}x_{162} - x_{161}^{2} - 2x_{162} + x_{162}^{2} + 2x_{163} \\*
&{} + x_{163}^{2} - 2x_{163}x_{165} - x_{164}^{2} - 2x_{165} + x_{165}^{2} + 2x_{169} + x_{169}^{2} - 2x_{169}x_{171} \\*
&{} - x_{170}^{2} - 2x_{171} + x_{171}^{2} + 2x_{172} + x_{172}^{2} - 2x_{172}x_{174} - x_{173}^{2} - 2x_{174} + x_{174}^{2} \\*
&{} + 2x_{178} + x_{178}^{2} - 2x_{178}x_{180} - x_{179}^{2} - 2x_{180} + x_{180}^{2} + 2x_{181} + x_{181}^{2} \\*
&{} - 2x_{181}x_{183} - x_{182}^{2} - 2x_{183} + x_{183}^{2} + 2x_{184} + x_{184}^{2} - 2x_{184}x_{186} - x_{185}^{2} \\*
&{} - 2x_{186} + x_{186}^{2} - 6x_{211}x_{213} - 4x_{213} + x_{214} \\[4pt]
% END_RESIDUAL 151
% BEGIN_RESIDUAL 152
R_{152} ={}& -2x_{97}x_{140} - 2x_{114}x_{167} - 2x_{121}x_{176} - 2x_{135}x_{216} + 2x_{139}x_{140} + 2x_{140} + 2x_{166}x_{167} \\*
&{} + 2x_{167} + 2x_{175}x_{176} + 2x_{176} + 6x_{211}x_{216} + x_{215} + x_{216} \\[4pt]
% END_RESIDUAL 152
% BEGIN_RESIDUAL 153
R_{153} ={}& 2 - 2x_{97} + x_{97}^{2} - 2x_{97}x_{139} + 2x_{97}x_{141} - 2x_{114} + x_{114}^{2} - 2x_{114}x_{166} + 2x_{114}x_{168} \\*
&{} + x_{121}^{2} - 2x_{121}x_{175} + 2x_{121}x_{177} + 2x_{139} + x_{139}^{2} - 2x_{139}x_{141} - x_{140}^{2} - 2x_{141} \\*
&{} + x_{141}^{2} + 2x_{166} + x_{166}^{2} - 2x_{166}x_{168} - x_{167}^{2} - 2x_{168} + x_{168}^{2} + x_{175}^{2} \\*
&{} - 2x_{175}x_{177} - x_{176}^{2} + x_{177}^{2} - 6x_{211}x_{215} - 3x_{215} + x_{216}
% END_RESIDUAL 153
\end{align*}
\endgroup

% END residuals.tex


\section{A reusable ordinate-normalization construction}\label{app:ordinate}
% BEGIN ordinate_lemma.tex
This appendix records a separate compression result for the retained
universal compiler. It is not offered as a better fixed Con(ZFC)
equation than the compact forms in Section~\ref{sec:quotient}.
The complete source and ordinary proof are under
\file{additional/ordinate_normalization/}.

Use the parent notation $\Delta=(a+2)^2-1$, with $c,i,T>0$ and
integer rank $R$, and put
\[
 I=ic^2,\quad S=\Delta I,\quad Q=S^2,\quad D=1+IS,
 \quad u=cT,\quad b=c+RD,\quad p=Du^2.
\]
For a new positive coordinate $z$, define
\[
 q_z=Qz,\qquad W=1+(q_z-z)(q_z-2),\qquad
 A=p+b^2-W,\qquad K=A^2-4pb^2.
\]
If $P_5$ is the product of the retained first, main, input, index and
transport factors, the new equation is
\[
 P_5(K+1)-1=0.
\]
It preserves the ordinary input of the universal compiler, with 17
positive witnesses, 91 operations and exact degree 269. The matched
older root-elimination construction has 92 operations and degree 325.
This universal count separates the ordinary input from the witnesses.

The mechanism is to impose $y=Qz-1$ before eliminating the strong
square root $f$. Indeed
\[
 QV^2-(Q-1)y^2-1=Q(V^2-W).
\]
The nonzero factor $Q$ cancels before taking the quadratic norm.
This avoids squaring a larger coefficient in the elimination step.

For soundness, the positive domain gives $Q>1$, $D>c$, $u>0$ and
$b\ne0$, even before a bound on $R$ is available. At a zero, the
integer unit $K+1$ is 1 or 2 modulo 4, so it is $+1$; hence
$P_5=1$ and $K=0$. The proof-only rational value
\[
 f=\frac{A}{2ub}
\]
satisfies $f^2=D$. A rational square root of an integer is integral.
Set $V=uf-b$ and $y=Qz-1>0$. Then $V^2=W$, the old auxiliary norm
is one, and the normalized strong norm is one.

The inherited source-specific signed-quotient theorem is needed to
recover $f>0$. If $f<0$, replacing $(f,T)$ by $(-f,-T)$ leaves all
old products unchanged and produces the forbidden positive-$f$,
negative-$T$ zero. This is an explicitly retained theorem for the valid
compiler recipe; it is not asserted for arbitrary numerical ports.

For completeness the canonical parent uses $R\equiv3\pmod4$ and
$y_* =\psi_S(R)$. Equation~\eqref{eq:oddres} gives $Q\mid y_*+1$,
so $z=(y_*+1)/Q>0$ is integral. The canonical strong and auxiliary
completion then makes $K=0$. This preserves accepted-input existence,
not every old coordinate tuple.

The degree calculation uses $\deg P_5=81$, $\deg Q=46$,
$\deg W=94$, $\deg A=94$ and $\deg K=188$.
The full nonzero leading form and an attained coefficient
$32\,\mathrm{Bm1}^{162}$ are supplied. Finite source and algebra
checks accompany the ordinary proof and its stated parent interface.

% END ordinate_lemma.tex


\clearpage
\begin{thebibliography}{9}
\small\raggedright\setlength{\itemsep}{3pt}
\bibitem{repo} V. Reshetnikov and contributors, \emph{ProveIt}, comparison
commit \file{536c08e7913e59bd42e8fca66e89944c920884ad},
\url{https://github.com/VladimirReshetnikov/ProveIt}.
The source ledger gives exact paths, retained notices, and hashes.
\bibitem{factorialfree} \emph{Removing factorial authenticity from
Diophantine consistency}, incoming archive
\file{ConZFC_Factorial_Free_Binomial_CRT_2026-10-25.zip}, at that snapshot.
\bibitem{coprime} \emph{Coprime product certificates and positive Pell
centers}, package \file{142-coprime-products-positive-centers},
at that snapshot.
\bibitem{linearcrt} \emph{Linear CRT}, incoming archive
\file{ConZFC_Linear_CRT_2026-10-25.zip}, at that snapshot.
\bibitem{pooled} \emph{Pooled Powers}, incoming archive
\file{ConZFC_Pooled_Powers_2026-10-25.zip}, at that snapshot.
\bibitem{projection} \emph{A direct quartic for Con(ZFC) in 234 integer
variables}, incoming archive
\file{ConZFC_Quadratic_Projection_234_2026-10-25.zip}, at that snapshot.
\bibitem{cofinal} \emph{Cofinal witnesses and centered carriers for
Diophantine consistency}, incoming archive
\file{ConZFC_Cofinal_Witnesses_Centered_Carriers_2026-10-25.zip},
at that snapshot. Its retained sources include the earlier rational
adjacent-Pell integrality argument, used and reproved here.
\bibitem{bayer} J. Bayer, M. David, M. Hassler, Y. Matiyasevich, and
D. Schleicher, \emph{Diophantine Equations over $\mathbb Z$: Universal
Bounds and Parallel Formalization}, arXiv:2506.20909,
\url{https://arxiv.org/abs/2506.20909}. Contextual background; the
direct construction does not invoke a universal polynomial as a black box.
\end{thebibliography}
\end{document}
